\section{D. The arithmetic action, its Kummer coordinates, and the two CM factors}
\label{sec:D}

We use $G_k=\operatorname{Gal}(\overline k/k)$, $K=\mathbb Q^{\rm cyc}=\mathbb Q^{\rm ab}$, and the cyclotomic character $\chi:G_{\mathbb Q}\to\widehat{\mathbb Z}^{\times}$. Compatible roots of unity identify $T\mathbb G_m$ with $\widehat{\mathbb Z}(1)$. A coordinate in this Tate module depends on that choice. The group in scans 156--170 is the arithmetic $G_{\mathbb Q}$; the larger loop-automorphism groups enter only in the explicitly designated comparisons below. The relevant scans 157,166--170 have been inspected as images, independently of their transcription. A positive-weight calculation always refers to the kernel of the multiplier and to its lower-central filtration.

\subsection{The homomorphism of scan 157 and the rational-prime questions}

\begin{statement}{D.1}{Proved}{The actual affine Kummer homomorphism}
Choose compatible roots $a^{1/n}$ of $a=-1/2$. The homomorphism of \S51.15, (326), (329)--(335), scans 156--160, is, after a choice of origin in its torsor,
\[
 \rho_a:G_{\mathbb Q}\longrightarrow
 \widehat{\mathbb Z}^{\times}\ltimes\widehat{\mathbb Z}(1),
 \qquad g\longmapsto(\chi(g),\kappa_a(g)),\qquad
 g(a^{1/n})/a^{1/n}=\zeta_n^{\kappa_a(g)\bmod n}.
\]
Its image is the closed subgroup of index two
\[
 \{(u,v):(-1)^v=\chi_{-8}(u)\},                         \tag{D.1}
\]
where $\chi_{-8}$ is the quadratic character of $\mathbb Q(\sqrt{-2})$, with values $+1$ at $1,3\pmod8$ and $-1$ at $5,7\pmod8$. Changing the root system translates the origin and changes $\kappa_a$ by $(\chi-1)b$. On $G_K$ its image is exactly $2\widehat{\mathbb Z}(1)$.
\end{statement}
\begin{proof}
The cocycle identity is obtained by applying $g$ to $h(a^{1/n})$:
$\kappa_a(gh)=\kappa_a(g)+\chi(g)\kappa_a(h)$. This is precisely the multiplication in the displayed affine group. In the homogeneous Fermat model of the manuscript, the point $1/2$ has coordinates $(-1/2,-1/2,1)$; its fibre is the diagonal pair of root torsors of $-1/2$. Dividing the two coordinates removes the diagonal ambiguity and gives the rank-one torsor of (329)--(335). Thus this construction supplies the stated morphism, rather than merely a group of the same abstract form. The unread signs in (332), scan 159, do not change this argument: the homogeneous equation and the specified point determine the two $n$th powers to be $-1/2$.

The square root $\sqrt{-2}/2$ lies in $K$, and hence the parity of $\kappa_a(g)$ is exactly the quadratic action on $\sqrt{-2}$. On $G_K$ the roots of $-1$ are fixed, so $\kappa_a=-\kappa_2$. D.2 below proves that this has image $2\widehat{\mathbb Z}(1)$. Finally $\chi$ is onto. Every fibre over $u$ is a coset of that translation image; the parity constraint therefore describes the entire fibre, proving (D.1).
\end{proof}

\begin{statement}{D.2}{Proved}{Exact joint images and independence for distinct rational primes}
For distinct positive rational primes $p_1,\ldots,p_r$, the continuous homomorphism
\[
 (\kappa_{p_1},\ldots,\kappa_{p_r}):G_K\longrightarrow
 \widehat{\mathbb Z}(1)^r
\]
has image $2\widehat{\mathbb Z}(1)^r$, of index $2^r$. Consequently every odd-primary projection is onto, the $2$-primary image is $2\mathbb Z_2^r$, and the characters are linearly independent over $\widehat{\mathbb Z}$. On $G_{\mathbb Q}$ the full affine image consists exactly of the pairs $(u,(v_i))$ satisfying
$(-1)^{v_i}=\chi_{\mathbb Q(\sqrt{p_i})}(u)$ for every $i$.
These assertions answer the individual and simultaneous questions of scans 168--169; integral surjectivity as written there is false.
\end{statement}
\begin{proof}
Let $a=\prod p_i^{e_i}>0$ and $n\geq1$. We first compute the exact relation
\[
 a\in K^{\times n}\quad\Longleftrightarrow\quad n\mid2e_i\quad\hbox{for every }i.       \tag{D.2}
\]
Choose the positive real root $\alpha=a^{1/n}$. Any other root differs by an element of $K$, so membership in $K$ is independent of this choice. A finite subfield $\mathbb Q(\alpha)$ of the abelian extension $K/\mathbb Q$ is Galois and totally real. Every conjugate of $\alpha$ is a real root of $X^n-a$, of which there are at most two. Thus $\alpha^2\in\mathbb Q$ and valuation at $p_i$ gives $n\mid2e_i$. Conversely this divisibility makes $\alpha$ a product of rational numbers and square roots of rational primes. Such square roots are cyclotomic: $\sqrt2=\zeta_8+\zeta_8^{-1}$, and the quadratic Gauss sum for an odd prime has square $(-1)^{(p-1)/2}p$. Multiplication by $i$ removes the latter sign. This proves (D.2).

For completeness, the finite Kummer pairing used here needs no rank estimate. Let $B$ be the subgroup generated by the classes of the $p_i$ in $K^\times/K^{\times n}$ and let $L=K(p_i^{1/n}:1\leq i\leq r)$. The pairing
$\operatorname{Gal}(L/K)\times B\to\mu_n$, $(g,a)\mapsto g(a^{1/n})/a^{1/n}$,
is nondegenerate on both sides: a Galois element fixing every root is trivial, and a root fixed by every element belongs to $K$. Both groups have exponent dividing $n$; injections into their finite character groups therefore have equal orders and are isomorphisms. By (D.2) the annihilator in $(\mathbb Z/n)^r$ of the relation subgroup is exactly $2(\mathbb Z/n)^r$. These are all finite images. Compactness gives the inverse-limit image stated above. A linear relation vanishes on $2\widehat{\mathbb Z}^r$ only when all coefficients vanish, since multiplication by two on $\widehat{\mathbb Z}$ is injective. Parity and surjectivity of $\chi$ give the last assertion just as in D.1.

The program \texttt{d\_kummer\_soule.py} enumerates these annihilators for every $n\leq24$ and records the indices for $r=1,2,3,4$. This checks, in particular, the finite obstruction to the proposed surjectivity.
\end{proof}

\begin{statement}{D.3}{Proved}{General openness and two corrections to the manuscript}
Let $k$ be a field finitely generated over $\mathbb Q$ and $A\subset k^\times$ a finitely generated subgroup of torsion-free rank $r$. The Kummer homomorphism on $G_{k(\mu_\infty)}$ has open image in $\operatorname{Hom}(A,\widehat{\mathbb Z}(1))\simeq\widehat{\mathbb Z}(1)^r$. This proves the openness requested under (355), scans 166--167, also for normal connected schemes essentially of finite type over $\mathbb Q$ and their groups of units. The stronger surjectivity proposed in scan 167 is false even for the full unit group of a normal arithmetic scheme of dimension two. The criterion for cyclotomic surjectivity proposed in scan 163 is also false.
\end{statement}
\begin{proof}
Torsion in $A$ contributes no homomorphism to the torsion-free Tate module, so choose a free complement. Write $H=\operatorname{Gal}(k(\mu_\infty)/k)$, an open subgroup of $\widehat{\mathbb Z}^{\times}$: the algebraic closure of $\mathbb Q$ inside $k$ is a number field. Choose an integer $m$ such that $H$ contains all units congruent to one modulo $m$. There is $u\in H$ with $u-1=dv$, where $d=\operatorname{lcm}(2,m)$ and $v\in\widehat{\mathbb Z}^{\times}$: take $u_p=1+p^{v_p(m)}$ at $p\mid m$, $u_p=2$ at odd $p\nmid m$, and $u_2=3$ when $2\nmid m$.

Suppose $\alpha^n=a\in A$ with $\alpha\in k(\mu_\infty)$. Put $\tau(\alpha)/\alpha=\zeta_n^{c_\tau}$. Since $H$ is abelian, comparison of its two cocycle products gives
$(u-1)c_\tau=(\chi(\tau)-1)c_\sigma$ for an element $\sigma$ with character $u$. Setting $b=v^{-1}c_\sigma$ modulo $n$ shows that $\alpha^d\zeta_n^{-b}$ is fixed by $H$. Thus
\[
 a^d\in k^{\times n}.                                      \tag{D.3}
\]
Choose a normal finitely generated $\mathbb Z$-algebra $R$ with fraction field $k$ containing $A$ as units. Its unit group is finitely generated. One proof is to take a normal projective compactification over $\operatorname{Spec}\mathbb Z$: a unit has divisor supported on the finitely many codimension-one boundary components, while the kernel of this divisor map is the unit group of the ring of global functions of the compactification. That ring is finite over $\mathbb Z$; its units are finitely generated by the number-field unit theorem. Normalization is finite here, so this compactification argument applies to $R$.

If $b^n=a^d$ in $k$, both $b$ and $b^{-1}$ are integral over $R$, hence belong to $R$. In the free lattice $U=R^\times/\mu(R)$, the saturation of $A$ modulo $A$ is finite; let $e$ annihilate it. Equation (D.3) then implies $de\,a\in nA$ in additive lattice notation. The kernel of $A/nA\to k(\mu_\infty)^\times/k(\mu_\infty)^{\times n}$ has order at most $(de)^r$. The perfect finite pairing in D.2 therefore bounds all finite indices of the Kummer image by this same number. A closed subgroup of $\widehat{\mathbb Z}^r$ with uniformly bounded finite indices has finite index: any collection of more distinct cosets is separated in one finite quotient. This proves openness.

For a normal connected scheme, the Kummer covers of units are finite etale and the generic-point map is onto their monodromy groups; the preceding field computation consequently gives the same image. For the counterexample take $S_1=\operatorname{Spec}\mathbb Z[1/2][t]$ and $A=\mathcal O(S_1)^\times=\{\pm2^j\}$. Its generic fibre has the Kummer image $2\widehat{\mathbb Z}$ of D.2. Moreover $A$ is saturated in $\mathbb Q[t]^\times$ modulo its torsion, so the proposed saturated variant fails too. Finally $k=\mathbb Q(\sqrt5)$ has only the roots of unity $\pm1$, but its cyclotomic image has index two because $k\subset\mathbb Q(\mu_5)$.
\end{proof}

\begin{statement}{D.4}{Proved}{Exact local inertia images}
Embed $G_{\mathbb Q_\ell}$ into $G_{\mathbb Q}$ and put $I_\ell^+=I_\ell\cap\ker\chi$. For a positive rational prime $p$, the image of $\kappa_{p,\ell}$ on $I_\ell^+$ is $\ell^{e_{p,\ell}}\mathbb Z_\ell(1)$, where
\[
 e_{p,\ell}=\begin{cases}
 0&\ell>2, p=\ell,\\
 v_\ell(p^{\ell-1}-1)-1&\ell>2, p\ne\ell,\\
 1&\ell=p=2,\\
 v_2(p^2-1)-2&\ell=2, p\text{ odd}.
 \end{cases}                                               \tag{D.4}
\]
Every image in scan 170 is therefore open, and the formula determines exactly which are onto. The program \texttt{d\_local\_cm.py} records every pair $p,\ell\leq50$; for example the image for $(p,\ell)=(3,11)$ has index $11$, and for $(19,7)$ index $49$.
\end{statement}
\begin{proof}
The fixed field of $I_\ell^+$ is $L=\mathbb Q_\ell^{\rm ur}(\mu_{\ell^\infty})$, abelian over $\mathbb Q_\ell$. For odd $\ell$, choose a cyclotomic multiplier with $u-1\in\mathbb Z_\ell^\times$. The cocycle argument in D.3, now with $d=1$, proves that an $\ell^k$th root of $a\in\mathbb Q_\ell^\times$ belongs to $L$ precisely when $a\in\mathbb Q_\ell^{\times\ell^k}$. The decomposition
$\mathbb Q_\ell^\times=\ell^{\mathbb Z}\times\mu_{\ell-1}\times(1+\ell\mathbb Z_\ell)$ and the $\ell$-adic logarithm give the first two cases.

For $\ell=2$ use multiplier $-1$. The same argument proves that $a$ having a $2^k$th root in $L$ implies $a^2\in\mathbb Q_2^{\times2^k}$, equivalently $a=\pm b^{2^{k-1}}$ for $b\in\mathbb Q_2^\times$. Conversely every quadratic extension of $\mathbb Q_2$ lies in $L$: the square classes are generated by $-1,2,5$, the first two have cyclotomic square roots and $\mathbb Q_2(\sqrt5)$ is unramified. Thus $\pm b^{2^{k-1}}$ indeed has a $2^k$th root in $L$. For $a=2$ the largest such $k$ is one. For odd $p$, write $p=\epsilon u$, $\epsilon\in\{\pm1\}$, $u\in1+4\mathbb Z_2$; logarithms give the largest $k$ as $v_2(u-1)-1=v_2(p^2-1)-2$. Finite Kummer duality converts these root heights into the displayed closed images.
\end{proof}

\subsection{The universal Fermat cocycle and its actual comparison maps}

\begin{statement}{D.5}{Known}{Full profinite universal measure, and the precise correction to Lead 1}
Put $\Lambda=\widehat{\mathbb Z}[[\widehat{\mathbb Z}^2]]$. The compatible projective Fermat homology of Otsubo's Theorem 6.1 is a free cyclic $\Lambda$-module. For its specified generator $c$, write $g(c)=F_g c$; then
$F_{gh}=g(F_h)F_g$. On $G_K$ this is a homomorphism into $\Lambda^\times$. The comparison with Anderson's relative-path measure $B_g$ is the identity
\[
 (1-s^r)(1-t^m)F_g=(1-s^{\chi(g)r})(1-t^{\chi(g)m})B_g.       \tag{D.5}
\]
Here $r,m$ are integers and $s,t$ are the two deck generators. This is Otsubo's (7.2). In particular $F_g=B_g$ on $G_K$. Its $\ell$-power specialization is related to Ihara's power series by the explicit derived-module basis comparison below.

The morphisms relating this invariant to D.1 are the exact quotient
\[
 1\longrightarrow\widehat\pi'/\widehat\pi''
 \longrightarrow\widehat\pi/\widehat\pi''
 \longrightarrow\widehat\pi^{\rm ab}\longrightarrow1
                                                               \tag{D.6}
\]
and the root torsors obtained by evaluating the abelian covering at $1/2$. For rational primes the additional exact comparison is
\[
 p=\prod_{a=1}^{p-1}(1-\zeta_p^a),\qquad
 \kappa_p\big|_{G_K}=\sum_{a=1}^{p-1}\kappa_{1-\zeta_p^a}\big|_{G_K}.       \tag{D.7}
\]
Lead1's proposed Kummer comparison is made exact by (D.5)--(D.7) and the coefficient extraction in D.23 below. In particular the rank-one affine coordinate is obtained from the full universal measure by a specified continuous map. This comparison retains the factor two in D.2 and the prime-to-$\ell$ deck components that a pure pro-$\ell$ series discards.
\end{statement}
\begin{proof}
We use \cite[Theorem 6.1, Corollary 6.2, equation(7.2)]{Otsubo2016}, whose full text was opened. The non-zero-divisor assertion in Corollary 6.2 permits cancellation in (D.5); when $\chi=1$ the two measures coincide. The maps to $\ell$-power deck quotients and $\mathbb Z_\ell$ coefficients are continuous quotient maps of completed group rings; hence this comparison is profinite before any pro-$\ell$ specialization. Part B constructs the derived-module comparison in (D.6); evaluation of its abelian quotient at the displayed point gives D.1 by the explicit root equations. Finally (D.7) is $\Phi_p(1)=p$ and functoriality of roots under multiplication. Applying valuations at the distinct primes to its rational norm products, followed by the exact pairing of D.2, proves independence through these cyclotomic elements. The possible quadratic relations, including the factor two, are retained; an assertion of unrestricted integral independence of all circular units would not give the required exact image.

Otsubo's Remark 8.2 identifies the original results precisely as Ihara1986, Theorem 7 for prime-power level, and Anderson1989, Theorem 6 in general. Anderson's complete original has now been retrieved from the official GDZ scan and checked. Its \S\S1.3--1.4 use the full ring
$\widehat{\mathbb Z}[[\widehat{\mathbb Z}(1)^2]]$;
Proposition 5.5.1, p.88, constructs the all-degree torsion comparison from relative interval homology and states
$e_N(B_\sigma\,\sigma(v))=\sigma(e_N(v))$.
Theorem 6, pp.118--124, then interpolates Gauss sums using a specified Frobenius lift and the hyperadelic gamma function; it is not by itself the torsion comparison.

The original comparison to Ihara is \cite[Proposition 13.3.2 and \S\S13.4--13.6]{dAnderson1989}.
For odd $\ell$ it constructs a Galois-equivariant map
$\theta:\pi_\ell'/\pi_\ell''\simeq
\varprojlim_n T_\ell J_{\ell^n}$ and a basis-change unit
$Q(u,v)$, with exact formula
\[
\begin{aligned}
 F^{\rm Ih}_\sigma(u,v)&=S_\sigma(u)S_\sigma(v)
 B_{\sigma,\ell}(u,v)
 \frac{Q((1+u)^{\chi_\ell(\sigma)}-1,
          (1+v)^{\chi_\ell(\sigma)}-1)}{Q(u,v)},\\
 S_\sigma(T)&=\frac{(1+T)^{\chi_\ell(\sigma)}-1}{T}.
\end{aligned}
\]
The map $\theta$ sends a deck transformation to its translation
on the compatible division points of the divisor
$(\delta_n)-(\infty_n)$, where
$\delta_n=(t_n,-t_n(1-t^{-1})_n)$ in \S13.5.
That section records the plus-model twist for $\ell=2$;
RemarkII in \S13.3 proves $F^{\rm Ih}_\sigma=B_{\sigma,\ell}$
on the cyclotomic kernel for every $\ell$.
These maps specify both the basis coboundary and the arithmetic twist.

The original Ihara1986 and Coleman1989 proofs remain to be obtained; the separate coefficient formula used below was checked in Ihara--Kaneko--Yukinari's primary paper and now also in Anderson's original. The retrieved original \cite[Supplement]{dIhara2000} corrects Ihara1986, Theorem 7 by replacing the Jacobi indices $(a,b,c)$ with $(-a,-b,-c)$, and corrects the signs in Theorem 10. It also instructs replacing Anderson1989, p.125 line11's group-ring identification by its inverse, and records the inverse deck convention on p.85. These corrections explain why a convention must be attached to the original theorem numbers; D.6 uses the explicitly checked Otsubo convention.
\end{proof}

\begin{statement}{D.6}{Known}{Frobenius specialization with nonsplit primes and the plus model}
Let $p\nmid n$, $f=\operatorname{ord}_n(p)$, $q=p^f$, and let $\theta:\mathbb F_q^\times\to\mu_n$ be fixed. For $a,b,a+b\ne0$ modulo $n$, put
\[
 j_n^{a,b}=-\sum_{t\in\mathbb F_q}\theta(t)^a\theta(1-t)^b,
\]
with nontrivial characters extended by zero. Otsubo's Theorem 8.1 identifies this number with the character specialization of the product of $F_{\operatorname{Frob}_p}$ along the orbit
$(a,b),(pa,pb),\ldots,(p^{f-1}a,p^{f-1}b)$. For $x^n+y^n=z^n$ the sum of these $j_n^{a,b}$ is the trace of geometric Frobenius on $H^1$ over $\mathbb F_q$. For the manuscript's $x_0^n+x_1^n+x_\infty^n=0$ it is instead the sum of
\[
 j_{n,+}^{a,b}=\theta(-1)^{a+b}j_n^{a,b}.                    \tag{D.8}
\]
\end{statement}
\begin{proof}
The cited theorem \cite[Theorem 8.1]{Otsubo2016} is formulated for an arbitrary residue degree; its product, rather than a single value at $p$ on a non-fixed character, is essential. The comparison with $H^1$ uses the ordinary arithmetic Frobenius action on the Tate module and its dual geometric action on cohomology; these have the same characteristic polynomial in this convention. The change of models is defined over $\mathbb Q(\epsilon_n)$ by
$(x_0:x_1:x_\infty)\mapsto(x_0:x_1:\epsilon_n x_\infty)$, $\epsilon_n^n=-1$. Its Frobenius discrepancy is $\epsilon_n^{q-1}=(-1)^{(q-1)/n}$. Projectively this multiplies the first two coordinates by the same element. The $(a,b)$ character is therefore multiplied by $\theta(-1)^{a+b}$, which proves (D.8). This retains the arithmetic descent of scan 171, including even $n$.
\end{proof}

\begin{statement}{D.7}{Proved}{Complete finite-field computation in the requested range}
The files \texttt{d\_jacobi.py}, \texttt{d\_jacobi\_field.cpp}, \texttt{d\_jacobi.out} and \texttt{d\_jacobi.json} compute every nondegenerate character value for every prime $p\leq50$ and $3\leq n\leq12$ with $p\nmid n$, over the exact field $\mathbb F_{p^{\operatorname{ord}_n(p)}}$. Every value satisfies $j\overline j=q$. Character sums and independent point counts agree for both models. Levels 1 and 2 have zero projective $H^1$, and the output separately retains the bad-reduction pairs $p\mid n$.
\end{statement}
\begin{proof}
For non-semiprimitive pairs the code finds a monic polynomial for which multiplication by its residue class has order $q-1$. This certifies that all nonzero elements of the quotient ring form a cyclic group, hence that the quotient is the required finite field. It enumerates the cyclotomic-number array
$C_{ij}=\#\{t\ne0,1:\theta(t)=\zeta_n^i,\theta(1-t)=\zeta_n^j\}$.
Discrete Fourier summation gives each $j$ as an integral polynomial modulo $\Phi_n$. All norm identities are checked in this exact quotient ring. Independently, the standard curve has
$n^2C_{00}+2n+n\,[\theta(-1)=1]$ points, and the plus curve has
$n^2C_{ee}+3n\,[e=0]$, where $\theta(-1)=\zeta_n^e$.

For the larger semiprimitive fields, $p^h\equiv-1\pmod n$, $q=p^{2h}$, the code uses and checks the exact identity $j=-p^h$. Indeed the characters are trivial on $\mathbb F_{p^h}^\times$. In the quadratic extension, summing the additive character over trace fibres gives the Gauss sum $G(\theta^a)=p^h\theta^a(\tau)$ for a nonzero trace-zero element $\tau$; the quotient $G(\theta^a)G(\theta^b)/G(\theta^{a+b})$ is $p^h$. This proves the shortcut. The smaller semiprimitive cases are also enumerated directly. The largest direct field is $\mathbb F_{47^5}$; it was actually enumerated. For an explicit arithmetic-twist check, $(p,n)=(43,6)$ gives standard trace $32$ and plus-model trace $-64$, with respective point counts $12$ and $108$. These computations verify the finite-field side of D.6 and its character conventions; they do not claim a numerical oracle evaluating an arbitrary abstract Galois element independently of the theorem.
\end{proof}

\subsection{Circular Kummer characters, regularity, and the first irregular prime}

\begin{statement}{D.8}{Known}{The checked Ihara--Anderson--Coleman coefficient formula}
For odd $\ell$, let $U=\log(1+u)$, $V=\log(1+v)$, $W=-U-V$. On $G_{\mathbb Q(\mu_{\ell^\infty})}$ the universal series has the expansion
\[
 F_g(u,v)=\exp\left(\sum_{m\geq3\ {\rm odd}}
 \frac{\kappa_m(g)}{(1-\ell^{m-1})m!}
 (U^m+V^m+W^m)\right),                                  \tag{D.9}
\]
where $\kappa_m$ is the compatible Kummer character of
$\epsilon_n^{(m)}=\prod_{a\in(\mathbb Z/\ell^n)^\times}(\zeta_{\ell^n}^a-1)^{\langle a^{m-1}\rangle_n}$.
The integral factor $1-\ell^{m-1}$ is a unit. Vandiver at $\ell$ is equivalent to surjectivity of all these $\kappa_m$ for $m=3,5,\ldots,\ell-2$.
\end{statement}
\begin{proof}
This is the reduction of \cite[Theorems A$_2$ and B, equations(9)--(11)]{dIKY1987} to the root convention of D.1. Their character is defined by the same quotient of a chosen $\ell^n$th root and its Galois translate, so no unnamed rescaling is involved. Their Theorem B identifies $\beta_m=(1-\ell^{m-1})^{-1}\kappa_m$, and Theorem A$_2$ gives the exponential with $\beta_m/m!$. The paragraph immediately following Theorem B proves the stated Vandiver equivalence. Both pages were inspected visually. Reducing modulo $\ell$ permits replacing $\zeta^a-1$ by $1-\zeta^a$: the resulting sign is an $\ell$th power. These are the exact characters computed next.

Anderson's original \S11.6, Proposition 11.6.1 and formula(13.7.3),
pp.117--118 and 130, give an independent normalization check.
His $\beta_m$ sums $i^{m-1}$ times the Kummer cocycle of
$1-\zeta_{\ell^n}^{-i}$ over every $0<i<\ell^n$.
On the cyclotomic kernel this is the cocycle of
$\zeta_{\ell^n}^{i}-1$, since their ratio is
$\zeta_{\ell^n}^{-i}$ and all of its $\ell$-power roots are fixed.
Splitting the sum into $\ell\nmid i$ and $i=\ell j$ gives
$\beta_m=\kappa_m+\ell^{m-1}\beta_m$.
Thus his exponential coefficient is exactly the one in (D.9).
\end{proof}

\begin{statement}{D.9}{Known}{An exact obstruction module, not an irregular-prime guess}
Let $\mathcal F_\ell\subset\mathbb Z_\ell[[u,v]]^\times$ be the group defined by the five conditions of Ichimura--Kaneko1989, \S1.2: congruence $F\equiv1\pmod{uvw}$, inversion $F\overline F=1$, symmetry in $u,v,w$, symmetry of $F*F$ in four variables with product of $1+$variables equal to one, and
\[
 \prod_{\zeta^\ell=1}F(\zeta(1+u)-1,\zeta(1+v)-1)
       =F((1+u)^\ell-1,(1+v)^\ell-1).
\]
Their Theorem 1 identifies equality $\operatorname{Im}F=\mathcal F_\ell$ with Vandiver at $\ell$. Their Theorem 2 computes the characteristic power series of the defect $(\mathcal F_\ell/\operatorname{Im}F)(-1)$ as that of
$\operatorname{Hom}_{\mathbb Z_\ell}(A_\infty^+,\mathbb Q_\ell/\mathbb Z_\ell)$ over the cyclotomic Iwasawa algebra. Thus the relevant obstruction is the plus class-group module, not merely divisibility of a Bernoulli numerator. At every regular odd prime this defect is zero. D.10 proves the same equality at every odd prime at most50, including 37.
\end{statement}
\begin{proof}
We use \cite[Theorems1--2]{dIchimuraKaneko1989}; the defining conditions and both theorem statements were checked on the original pages 314--315. Equation(D.9) identifies their $F$ and Kummer characters with ours. The plus part $A_\infty^+$ is the direct limit of the $\ell$-Sylow parts of the ideal class groups of $\mathbb Q(\mu_{\ell^n})$, with the authors' transition maps. The stated characteristic-series equality records exactly the size invariant proved there; it is not promoted to an isomorphism of modules. Regularity implies that the full class number at the first level is prime to $\ell$, hence its plus part is, and so Vandiver holds. For the bounded range, the certificates below establish the equivalent surjectivities directly.
\end{proof}

\begin{statement}{D.10}{Proved}{Soulé characters modulo every prime at most50}
For every odd prime $\ell\leq50$ and each odd $3\leq m\leq\ell-2$, the program \texttt{d\_kummer\_soule.py} exhibits an explicit nonzero value of $\kappa_m\bmod\ell$. Consequently all these individual characters are onto $\mathbb Z_\ell$, their distinct Teichmuller-weight components are independent, and Vandiver holds throughout this range. At $\ell=37$ all seventeen characters are certified by the single prime $q=5477$. The Bernoulli computation finds the irregular pair $(37,32)$ and no other irregular pair below 50; it causes no missing plus-part image here.
\end{statement}
\begin{proof}
For a prime $q\equiv1\pmod{\ell^2}$, a primitive root $g\bmod q$, and $r=g^{(q-1)/\ell}$, the code evaluates
$c_m(r)=\prod_{a=1}^{\ell-1}(1-r^a)^{a^{m-1}\bmod\ell}$ and finds the unique $k\bmod\ell$ such that
$c_m(r)^{(q-1)/\ell}=r^k$. A nonzero $k$ proves that the corresponding cyclotomic element is not an $\ell$th power. All choices, exponents and residues are retained in the JSON, so each witness is checked by modular exponentiation alone.

This is also a certificate on the infinite cyclotomic field, not only at its first layer. In $\mathbb Q(\zeta_\ell)^\times/\ell$ the class of $c_m$ has Teichmuller character $b\mapsto b^{1-m}$. The only order-$\ell$ subextension of the cyclotomic $\mathbb Z_\ell$-extension has Kummer class $\zeta_\ell$, of character $b\mapsto b$. These characters differ in the indicated range. Thus the nonzero class cannot disappear in that extension. The character takes a unit value in $\mathbb Z_\ell$, proving surjectivity; projection by the distinct idempotents of $\mathbb F_\ell^\times$ proves joint independence of these weight components. D.8 now proves Vandiver without inferring it from the irregularity test.

The witness primes, grouped by $\ell$, are
\[
\begin{array}{c|rrrrrrr}
\ell&5&7&11&13&17&19&23\\\hline
q&101&197,491&727&677,2029&3469&10831&12697
\end{array}
\]
\[
\begin{array}{c|rrrrrr}
\ell&29&31&37&41&43&47\\\hline
q&10093,13457&15377&5477&16811&22189&8837.
\end{array}
\]
For $\ell=3$ the specified weight list is empty; for $\ell=2$ there is likewise no weight in the requested range. No omitted character is hidden by these conventions.
\end{proof}

\subsection{Metabelian leading terms and the lifting issue}

\begin{statement}{D.11}{Proved}{Exact special-derivation images through weight 10}
Let $L$ be the free Lie algebra over $\mathbb Q$ on $x,y$, with the three peripheral elements $x,y,-x-y$. Modulo inner derivations normalize a positive-degree special derivation by $D(x)=0$, $D(y)=[f,y]$. For weight $n\geq2$, it is special exactly when some $g\in L_n$ satisfies
\[
 [f,y]=[g,x+y].                                           \tag{D.10}
\]
The metabelian quotient is
$L/L''=\mathbb Qx\oplus\mathbb Qy\oplus\mathbb Q[X,Y]c$, where $c=[x,y]$, $X=\operatorname{ad}x$, $Y=\operatorname{ad}y$; the last summand is abelian. The following exact rational ranks were computed, with source witnesses, by \texttt{d\_metabelian.py}:
\[
\begin{array}{c|rrrrrrrrrr}
n&1&2&3&4&5&6&7&8&9&10\\\hline
\text{special outer source}&0&0&1&0&3&0&6&4&13&12\\
\text{its metabelian image}&0&0&1&0&3&0&5&0&7&0\\
S_3\text{-invariant source}&0&0&1&0&1&0&2&1&4&2\\
\text{its metabelian image}&0&0&1&0&1&0&1&0&2&0
\end{array}                                               \tag{D.11}
\]
Thus all three weight-five directions survive. The extra invariant direction in weight seven dies. In weight nine two invariant directions survive. These are images of derivations of the free Lie algebra; special derivations defined only on $L/L''$ have dimension $n-2$ for every $n\geq3$, and are a larger calculation.
\end{statement}
\begin{proof}
Subtracting the inner derivation that conjugates $x$ gives the normalization. For $n\geq2$, the centralizer of $x$ has no component of degree $n$, so it is unique. The condition at the third peripheral element gives (D.10), with a harmless simultaneous sign determined by the convention for $g$. The free metabelian description follows by Jacobi: the commuting operators $X,Y$ generate $L'/L''$ from $c$, and its monomials are independent, for instance in the usual upper-triangular matrix realization over $\mathbb Q[X,Y]$.

The program uses the Lyndon basis in each degree and expands its brackets in the free associative algebra. It forms the exact matrix of $(f,g)\mapsto[f,y]-[g,x+y]$ and computes its rational nullspace. The invariant subspace is separately defined by the checked LM2 equations
\[
 f(x,y)+f(y,x)=0,\qquad
 [x,f(x,y)]-[x+y,f(-x-y,y)]=0.                             \tag{D.12}
\]
These follow by conjugating the normalized outer derivation by the transposition and the three-cycle of the three peripheral elements and then renormalizing $D(x)=0$. Thus the matrix computes the actual infinitesimal centralizer, rather than the two linear cycle equations. Projection of each Lyndon bracket to $\mathbb Q[X,Y]c$ gives the image ranks. Every pivot-image vector retains an exact preimage in the JSON.

In the metabelian algebra (D.10) reads $Yp=(X+Y)q$, hence $(X+Y)\mid p$. This gives the stated $n-2$ comparison dimension. At weight five the full image consists of $(X+Y)$ times all degree-two polynomials; the invariant image is spanned by
$X^3+2X^2Y+2XY^2+Y^3$.
At weight seven its invariant image is spanned by
$X^5+3X^4Y+5X^3Y^2+5X^2Y^3+3XY^4+Y^5$.
At weight nine one extra independent vector is $X^2Y^2(X+Y)^3$; the output gives the other vector and both preimages. The zero even-weight images, including weights eight and ten with nonzero sources, are computed by this projection, not assumed from parity.
\end{proof}

\begin{statement}{D.12}{Proved}{The reviewed pro-$\ell$ lifting assertion and its precise scope}
For a free pro-$\ell$ group on finitely many peripheral generators, every integral special derivation is the leading term of a special automorphism. For type $(0,3)$ and $\ell\geq5$, every $S_3$-invariant integral leading term lifts to an element of the centralizer of the geometric $S_3$ in the pro-$\ell$ outer special-automorphism group. Thus the image ranks in (D.11) are ranks of actual pro-$\ell$ special-automorphism images, and of their centralizer for $\ell\geq5$.
\end{statement}
\begin{proof}
Write $\Gamma_m$ for the closed lower-central series. Its associated graded is the free Lie algebra over $\mathbb Z_\ell$. Start with conjugators representing the given degree-$n$ Lie elements. The product-of-peripheral-generators equation holds modulo $\Gamma_{n+2}$ because its leading error is the special-derivation equation. At stage $m$ the error lies in $\Gamma_m/\Gamma_{m+1}$; changing the conjugators by elements of $\Gamma_{m-1}$ changes it by
$\sum_i[a_i,x_i]$. This map is onto $L_m$: iterated Jacobi rewrites every Lie monomial of degree $m$ as a sum $[x_i,b_i]$. In the two-generator normalization fixing $x$, the two available operators are $[\ ,y]$ and $[\ ,x+y]$; their sum still spans $[x,L_{m-1}]+[y,L_{m-1}]=L_m$. We can therefore cancel each error without altering the prescribed leading term.

The corrections converge in the free pro-$\ell$ group. Indeed every finite $\ell$-group quotient is nilpotent, so $\Gamma_m$ eventually lies in every open normal subgroup. The resulting endomorphism is the identity on the Frattini quotient and hence is onto by the pro-$\ell$ basis theorem; a finitely generated profinite group is Hopfian, so it is an automorphism. This independently proves the reviewed assertion.

For the invariant assertion, the group of IA automorphisms in this pro-$\ell$ setting is pro-$\ell$, as are the successive kernels. A finite group of order prime to $\ell$ acts on it. For a quotient by a stable closed normal subgroup, fixed elements lift to fixed elements: for finite $\ell$-groups this follows by induction on a central series and averaging the additive obstruction, and the profinite assertion follows by compactness. Apply this to the order-six group $S_3$ and successive filtration quotients. This proves the claimed lift. No such convergence argument applies to the full free profinite group: a finite nonnilpotent quotient need not kill any $\Gamma_m$.
\end{proof}

\begin{statement}{D.13}{Open}{The remaining full profinite image comparison in D(e)}
The maps from actual full profinite loop-preserving outer automorphisms, and from $\widehat{\rm GT}_0$, to their pro-$\ell$ leading terms land in the second and fourth rows of (D.11), respectively. D.12 proves the reverse inclusions for the pro-$\ell$ groups, not for those full profinite source groups. The full profinite realization of the extra weight-five metabelian directions, and of the extra invariant weight-nine direction, remains unproved here. Consequently D(e), as asked for the three full profinite groups, remains undecided in this return.

D.39 extends a new representative through the entire all-prime mixed Fermat homology tower, jointly with the retained degree-eight quotient. Its common quotient stabilizes at exponent eight. The exact two-adic system has cokernel killed by sixteen, and the half-average gives every odd-prime component. The resulting multiplier-one automorphism retains the normalized G6 and T3 actions, replaces the prior J3 and J12 actions by the specified new components, and retains its checked J48 seed at every later level. Both compact CM matrices and the fixed-fibre pair survive. The exact source kernel is the intersection of the closed L-commutator subgroup with D9K; the nonzero degree-sixteen word remains in it. The projection to the earlier finite joint group is strict, detected by $x^{144}$. D.45 proves that this recorded action has no full profinite lift. D.46--D.47 subsequently construct a replacement on their exact relative sources; its full profinite realization remains unresolved.

D.40 extends that action through the actual image in every finite
perfect quotient of the marked source. Their exact common quotient
with the retained tower is $C_6$, and the normalized conjugators pair
with the identity on all perfect coordinates. A further actual
$S_7$ quotient gives the fibre product over the common sign $C_2$.
The exact remaining full-source kernel maps onto $A_{10}$ under the
computed $S_{10}$ quotient. These maps retain both compact-cover and
fixed-fibre CM coordinates; D.45 excludes a full profinite lift of this recorded action.

D.41 then extends the action through all actual symmetric quotients
simultaneously. Their even image is the full product of the
alternating factors indexed by distinct marked kernels, and the
common quotient with D.40 is exactly the sign $C_2$.
Both projection kernels are computed. The remaining full-source
kernel maps onto $\operatorname{PSL}_2(\mathbb F_7)$ under the
specified $\operatorname{PGL}_2(\mathbb F_7)$ map and kills the
canonical perfect mod-seven coordinate. This retains the actual
CM data; D.45 excludes a full profinite lift of the recorded action.

D.42 extends this actual source action through the full relative
pronilpotent kernel of the specified projective-linear quotient.
Its maximal common marked quotient has base order $54$ and kernel
$\mathbb Z_3^{10}$. The nonzero common conjugator is retained by
an actual coprime centralizer lift, and every projection kernel is
computed. The remaining source kernel maps onto
$\operatorname{PSL}_2(\mathbb F_{11})$ under the specified
$\operatorname{PGL}_2(\mathbb F_{11})$ map. D.45 supplies a finite sphere quotient that excludes every full profinite lift of this recorded action.


D.43 extends the retained action simultaneously through all relative
pronilpotent completions whose marked finite base has nonabelian
simple even subgroup. The exact common quotient with the retained
tower is $C_{54}\times_{C_6}D_{\mathscr J}$, with its actual
source-kernel formula. One compatible centralizer lift supplies the
conjugator over the whole family. The exact remaining source kernel
maps onto $A_7\times A_7$ under an explicit quotient whose involution
exchanges the two simple factors. D.45 excludes every full profinite lift of this recorded action.

D.44 extends the same quotient action to the complete family of
anabelian even bases with their relative pronilpotent kernels.
The common quotient and both source kernels are exact, and an
actual $A_5^4\rtimes S_4$ marking detects residual kernel $A_5^4$.
D.45 then excludes a genuine full profinite lift of the entire
recorded normalized D.39 action and every D.40--D.44 extension
retaining it. Its sphere projection through $B/D_9F$ has a
nonzero degree-nine obstruction after every lower centralizer
and multiplier freedom is included. The elliptic projection
has an explicit corrected lift to $B/D_{10}B'$. D.46--D.47 subsequently supply a replacement realizing the prescribed
CM pair on their exact relative sources. The full source image
questions in D.13 and D.19 remain to be determined.

D.46 constructs a replacement on the entire relative pronilpotent
completion of the actual rank-97 subgroup $N$. Its full relative
pro-two component is an actual arithmetic commutator action with
identity sphere projection and compact elliptic scalar $\eta_*$;
every odd relative-prime component is the identity. The two actions
join through explicit coprime centralizer lifts. The resulting
$G_2,G_3$ actions determine the prescribed pair $T_4=I,W_3=[i]$,
and the reconstructed parameter has an actual $\widehat F'$ preimage.
Thus the replacement solves the entire rank-four relative pro-two
lifting step. Its exact remaining full-source kernel has image
$A_7$ under the displayed $S_7$ quotient. Its full profinite
realization remains in the lifting problem.

D.47 extends this replacement coherently over every Fermat level.
For \(M=\widehat N\cap\widehat F'\), it proves both inclusions
identifying the closed relative pro-prime source with the inverse
system over the actual \(L_m\). The common marked quotient is
\(\mathcal H_\infty\), and the simultaneous kernel is
\(\prod_pM^{(p)}\). The source kernel lies in \(M'\subset
\widehat F''\), so the action descends to the full metabelian
quotient, with every profinite deck coordinate retained. Its
coefficient-two action is the actual arithmetic \(h\) action;
every odd coefficient action is identity. The actual marked
cubic-four and quartic-three cover maps give \((-I,I)\), outside
the proved Galois joint graph. The Heisenberg quotient and full
Alexander-module argument also proves this in the outer action
space. All comparison kernels to D.46 are computed. The remaining
full-source kernel still has exact \(A_7\) image in the displayed
\(S_7\) quotient. The full GT0 and full loop-preserving
\(\widehat F\) lifting questions remain unresolved.

D.48 extends the entire D.47 action over the central double cover
\(\mathcal E=\widehat B/[\widehat N,\widehat B]\), jointly with
every Fermat level. The new subgroup \(M^+\) has index two in
\(M\); its relative pro-two source is unchanged. The comparison
kernel is the product of the explicitly identified odd-prime
kernels. The actual \(h\) action fixes the enlarged base, and
the coprime centralizer corrections retain the complete D.47
action and parameter. On the new finite mixed quotient
\(\mathcal J_3\), the exact old source-kernel image is
\(C_3^{96}\); the new source kernel has trivial image there.
Its full common quotient with D.47 is
\(\mathcal E\times_QG_3\), and the identity action on
\(\mathcal J_3\) agrees with the retained action on that whole
quotient. The remaining full-source kernel has exact
\(A_7\) image; full profinite realization remains unresolved.


D.49 extends this same action over the maximal marked quotient
fixed pointwise by the actual arithmetic \(h\), and over every
relative-prime source above that base. Its exact comparison
with D.48 retains both kernel layers, including an actual
\(A_5\) quotient of the base residual. The full congruence
quotient \(\operatorname{SL}_2(\widehat{\mathbb Z})/\{\pm I\}\),
with one global sign, factors through the identity-three
component; the new action is therefore pointwise identity
on this entire quotient. Its remaining source kernel has
exact image \([V_{3,2},A]\) in the specified mixed homology
quotient over \(G_3\). It contains an actual cyclic moving
subspace of dimension \(64\cdot3^{97}\); the dimension
of the full arithmetic span \(A\) is retained as unevaluated.
The earlier central iteration stabilizes by the exact
integral equality \([L,\widehat B]=L\).
The remaining mixed source and full profinite realization
questions stay active.


D.50 extends this same retained action through the entire
two-adic homology tower over the actual marked \(G_3\).
The exact residual image is the integral moving summand
\(U^-\), and every common quotient is computed from the
original source kernels. A regular-edge Fitting projection
and its explicit crossed homomorphism repair the initial
peripheral endomorphism to an automorphism; the cyclic
norm correction fixes the original cusp. Compatible
conjugators and invertible reduction modulo two give all
higher lifts and their continuous inverse limit. The
full parameter, CM pair, metabelian and congruence actions
are retained. The remaining source kernel has a nonzero
image containing \(W_1\wedge V\) in the next actual
class-two exponent-four quotient. Extending across that
remaining kernel and attaining the full profinite source
remain active steps.


D.51 retains this entire chosen homology action and both
full conjugators and extends them through the moving
part of the next actual class-two exponent-four quotient.
The exact old source image is \(Z=W\wedge V\); its
invariant and moving dimensions are respectively
\(m^2(q-1)/2\) and \(w\rho/2\). The actual common
quotient is \(\mathcal T_2/Z\). Characteristic descent,
exterior preservation of \(Z^A\), and Frattini generation
prove the new action on the joint source. Its projection
adds the whole moving central summand and retains every
previous parameter, CM, metabelian and congruence comparison.
The remaining exact image is \(Z^A\ne0\). In the separate
order-1080 verification, the complete corrected cusp word
has 4,159,444 letters. That example's fixed homology matrix
has intersection-pairing defect of rank 72, excluding its
extension to the unquotiented class-two peripheral source.
The arithmetic invariant obstruction and the full
profinite realization remain to be evaluated.


D.52 solves the entire invariant boundary equation in the
separate marked order-1080 source. Its 120-dimensional
space consists of allowable conjugator pairs. A retained
weight-60 parameter solves all thirteen actual quadratic
equations; the complete 6,352,056-letter corrected word,
both torsion relations and generation are verified.
The new homology action has deviation rank 96 over
\(\mathbb F_2\) and preserves the mod-two ribbon form.
It replaces the excluded finite representative while
retaining the actual invariant quotient. It also has a
complete two-adic homology continuation with the entire
new class-two action retained, exact common quotient and
both projection kernels. Its next exponent-eight residual
is twice the commutator module, witnessed by a nonzero
272-letter commutator square. The arithmetic \(G_3\)
correction and the full profinite realization remain to
be evaluated through their own actual source maps.


D.53 computes the full actual moving peripheral-pair space,
its exact homology image and kernel, and a positive arithmetic
dimension bound. Its integral lifts give actual doubled
automorphisms of the complete D.50 source. D.54 lifts every
such difference to the full joint class-two source, through
an explicit compatible square module. The fixed arithmetic
cusp defect is preserved literally, and the actual affine
quadratic map descends through every same-action ambiguity.
Its arithmetic zero locus and full profinite realization
remain to be evaluated. In the separate order-1080 source,
the full finite homology zero locus has exactly 864 elements;
all 262,144 parameters, the original equations, generation
and the complete long cusp words have been checked.


D.55 gives the exact full-group norm from the arithmetic
cusp obstruction to the compact orientation-bivector defect.
Its local image lies in the norm kernel, and its residual
quotient has explicit Tate-norm and compact-projection maps.
Every point of the reduced boundary zero set has invertible
homology, so generation follows there. The actual ordinary
relative-$H_2$ orientation comparison is an isomorphism.
The arithmetic zero set and residual norm kernel remain
unevaluated, with all maps retained for the next attack.


D.56 evaluates an actual irreducible $C_3^4$ quotient of
the arithmetic normal span and the unique parameter $(2,1,2,1)$.
The marked quotient has order 3888 and actual action
$\operatorname{Int}(t^6)$. Its exact $A_4$ comparison kernel
gives surjective integral homology and class-two boundary maps.
The smaller arithmetic correction space has dimension 40,
with invariant central target dimension 2560. Its actual
quadratic zero set and full profinite realization remain
unevaluated; every source and comparison kernel is retained.


D.57 now solves the complete ten-coordinate boundary equation
on the actual order-$3888$ arithmetic quotient, for both integral
rank-nine branches and all sixteen framed cases. The literal
genus-three cover and degree-eight elliptic map, exhaustive integral
conjugator search, separate plus-central alignment and augmentation
correction give the asserted automorphism of
$\widehat B/(R_h\cap M_*\cap L_{*,2})$.
Its entire $W_h$ action and prescribed normalized parameter are
retained. The exact remaining quotient kernels carry the full
$G_3$ invariant and full profinite realization questions onward.


D.58 retains the entire D.57 action and its conjugators while
restoring the full integral arithmetic homology. The actual
projective repair stays in
$\mathcal L=\ker(U^-\to U_*^-)$.
It then extends through the remaining moving class-two module,
with exact residual kernel
$\mathcal K_0=\ker(Z\to Z_*)^A$.
All $2560$ solved invariant coordinates are retained; the remaining
dimension is $m^2(|A|-1)/2-2560$.
The next actual invariant boundary and the full profinite realization
remain to be evaluated across this specified source map.


D.59 corrects the full homology on every active cusp-safe free
character orbit of the actual arithmetic source. Its invariant
central block has an injective full norm and compactification;
the exact unquotiented orientation identity therefore kills its
boundary. The residual kernel has dimension
$48w-2560-5\,308\,416\kappa$, where
$\kappa\geq830\,900\,448\,697\,714\,719\,719$.
It contains the specified $366\,080$-dimensional small-character
submodule. Its remaining boundary and full profinite realization
are the next actual-source steps.

D.60 now repairs twelve further central blocks on the same
arithmetic source. Its integral projective-cover construction
retains the complete D.57 action and the prescribed $\mathcal W_h$
projection, while changing the D.59 homology action on two
explicit character orbits. The full norm and compactification
are injective on the new $12\,288$-dimensional central block,
so its twenty-eight boundary coordinates vanish by the
unquotiented orientation identity. The explicit residual
small-character module has dimension $353\,792$ and
$612$ coinvariants; its actual remaining boundary and the
full profinite realization remain to be evaluated.

D.61 makes the exact integral homology replacement on the two
size-twelve character orbits. It retains D.57 and $\mathcal W_h$,
and its injective norm and compactification kill the actual
orientation boundary on seven previous self blocks and twelve
new cross blocks. The resulting lift retains $22\,016$ further
central coordinates and fifty coinvariants. Its specified
remaining module has dimension $331\,776$ and $562$ coinvariants.
At the D.61 stage the four actual involution-cover divided norms
leave twenty-six named self-wedge classes. D.62 evaluates their
arithmetic correction by the actual source words below.

D.62 cancels all twenty-six projective self-wedge classes by
actual source-involution and infinite-cusp correcting words.
The full cusp wrap is retained; its extra commutator is killed
by the specified projector. The literal comparison on $Y$ is
$\operatorname{Int}(b_{t,Y})\alpha_{61,Y}$, with unchanged
homology $T_{61}$ and retained D.57 and $\mathcal W_h$ maps.
The new lift retains $6656$ further central coordinates and
twenty-eight coinvariants. Its specified residual has dimension
$325\,120$ and $534$ coinvariants. The arithmetic boundary on
this residual and the full profinite source realization remain
the active original questions.

D.63 orthogonalizes the actual compact projective summands
through its explicit integral graph and lifts them inside the
retained source kernel. The actual peripheral derivation gives
$T_{63}=(1-e)T_{61}(1-e)+e$ and the exact lower action
$\operatorname{Int}(b_{t,Y})\beta_Y$, retaining the complete
D.57 and $\mathcal W_h$ maps and all old projector kernels.
The new lift adds $40\,960$ central coordinates and fifty-two
coinvariants. Its specified remaining kernel $\mathcal J_{63}$
has dimension $284\,160$ and $482$ coinvariants; its compact
norm has rank six and kernel dimension $476$.
The arithmetic boundary on this literal residual and the
full profinite source realization remain the original active
questions.

The successive residual descriptions above refer to their
numbered stages. D.64 extends the actual action through the
onto quotient $\Phi_{64}$ of $\mathcal J_{63}$ of dimension
$1930$, with fourteen coinvariants. Marked compact isometries
and actual source-involution words give
$T_{64}=T_{63}+A_1m_1+A_4m_4$ and its literal lower-action
comparison, retaining both D.57 branches, all sixteen framings,
the full $\mathcal W_h$ maps and source conjugators.
Its post-word moving correction uses the actual
$R_{50}\cap L_{*,2}$ source lifts before the final augmentation.
The exact remaining kernel $\mathcal J_{64}=\ker\Phi_{64}$
has specified dimension $282\,230$.
The fourteen uncovered compact-test coordinates retain
their test scope; the entire residual arithmetic boundary
and the full profinite source realization remain active.

D.65 retains the full second and third compact blocks by an actual
integral source correction. Its signed exterior norm divided by two
has rank six in each ten-dimensional test. Coinduction realizes the
full image, and the actual source-involution words kill both rank-four
kernels. The further quotient has dimension $5752$ and twelve
coinvariants, with total retained dimension $7682$ and twenty-six
coinvariants. Both D.57 branches, literal small conjugators and the
full $\mathcal W_h$ maps are retained. The lower action is exactly
$\operatorname{Int}(b_{t,Y})\beta_Y$ with its stated derivation and
all thirty-six cusp factors. The remaining specified kernel has
dimension $276\,478$ and maps onto the old compact norm kernel of
dimension $1918$. Its two remaining tested coordinates retain that
test scope. The residual source boundary and full profinite
realization continue to require their actual lifts.

D.66 fills the complete fourth compact comparison block by actual
quadratic source automorphisms. Its integral projective cover and full
quadratic source square realize both remaining compact directions,
retaining the literal small conjugators, both D.57 branches, all
$\mathcal W_h$ maps and the preceding three compact blocks. The total
retained compact quotient has dimension $9600$ and $28$ coinvariants.
The specified remaining source kernel is $J_{66}$, of dimension
$274\,560$. Its cusp-radical correction and the full profinite
realization continue to require their actual source maps.
\end{statement}
\begin{proof}[Established comparisons and attempted lifting]
Taking the maximal pro-$\ell$ quotient is characteristic, so it defines the displayed continuous homomorphisms of outer automorphism groups. Taking a leading term is injective on each filtration quotient by its definition; the peripheral and centralizer equations give (D.10)--(D.12). D.12 realizes every admissible leading vector after the change of source to the pro-$\ell$ group. A proposed full profinite lift would have to solve the same conjugacy and generation equations in every finite quotient, including nonnilpotent quotients. The lower-central correction sequence controls none of the residual errors in such a quotient. We also checked the weaker route of solving just the metabelian Lie equations: its larger dimension $n-2$ proves that this route discards further lifting conditions already at weight four. These two failed transfers are exact gaps, not counterexamples to a full profinite lifting theorem.

For the Galois source, (D.9) gives at most one rational direction in each odd weight and zero in even weights in this metabelian image. D.10 supplies nonzero directions through weight ten at $\ell=37$; E.3 supplies the rational general-prime comparison. The weight-five pro-$\ell$ image therefore distinguishes the whole special group from its $S_3$ centralizer and from Galois. D.20--D.21 below add the full profinite beta/gamma compatibility and prove equality of the rational metabelian graded images of Galois and full $\widehat{\rm GT}$ through weight ten. The pentagon used there is an additional equation on $\widehat{\rm GT}_0$. The exact missing lift test for the latter is D.22; the pro-$\ell$ construction alone supplies only the subfamily of its tests with nilpotent quotients. D.24 gives an explicit test of the general automorphism-lifting route and its peripheral limitation. D.25 constructs nonarithmetical finite GT shadows with generation verified. D.26 lifts all its selected actions through the complete abelian-kernel tower. D.27 proves that the twenty-three nonidentity selected actions are nevertheless nongenuine, and constructs an alternate action in the same extra CM pair that survives both characteristic class-two tests and extends over the full relative pro-three kernel. D.28--D.29 further extend the alternate homology class through a compatible mixed-prime quotient, the entire class-two pro-two tower, and every odd relative pro-prime kernel. D.30 excludes the two recorded later central representatives using a complete lower-centralizer calculation, then constructs a further central replacement. D.31 replaces that central representative again by an explicitly compared action on $\mathcal S_2^{(7)}$, whose kernel over $Q$ has class three and third subgroup $C_2^{1860}$; it retains the same homology and CM data and combines with every odd relative kernel and the mixed quotient. The canonical map to D.30's target has kernel $C_2^{1620}$, and the induced action differs from the earlier one already on $H^{(1)}$. D.32 then proves that this recorded degree-seven action has no peripheral lift even to the full relative pro-two source: its next central layer, all lower centralizer freedoms, and the multiplier are exhausted by an explicit separating functional. D.33 constructs a new degree-eight replacement on the exact quotient $\mathcal S_2^{(8)}$, with class-four kernel, third subgroup $C_2^{9105}$ and fourth subgroup $C_2^{1620}$. Its canonical projection to D.31 has kernel $C_2^{7245}$, and its induced action changes the recorded image of $s$ first in degree five. A restricted stable correction extends that new action over the entire class-two tower while retaining degree eight, all odd relative kernels and the mixed quotient. D.34 identifies the full relative pro-two source with the exact outer deck centralizer of the marked genus-one three-cusp cover, including the full multiplier pair and the cubic arithmetic twist. D.35 constructs actual odd-deck boundary lifts and proves that the whole compact CM torus is already realized by actual Galois action centralizing all deck maps. Its exact reduction map retains the finite-action fibre over the recorded D.33 action. D.45 proves that particular fibre empty by an actual source-inner transport; the second CM factor remains to be realized with a replacement action. D.36 supplies the second marked elliptic cover and the exact shared cyclic-twelve homology maps. Its arithmetic joint image has order six, with the cubic and tangent twists retained; it also computes the gauge-invariant discrepancy between the D.33 cover sign and its fixed-fibre $T_4$ coordinate. D.37 then constructs an actual generating shadow on the shared congruence/Fermat-twelve exponent-twelve homology quotient. Its exact pre-gauge $G_6$ conjugators agree with the retained D.33 components, while its actual normalized compact-cover matrices are $-I$ on cubic four-torsion and $I$ on quartic three-torsion, both evaluated through one cyclic-twelve cover. The two-cover arithmetic graph fails on this finite shadow. D.37 also computes both universal relative-prime quotients and the exact mixed kernel $(\mathbb Z/4)^{768}\oplus(\mathbb Z/3)^{864}$. D.38 constructs a new action on the exact fibre product of this Fermat quotient with D.33's degree-eight quotient. It proves both kernel maps and the common quotient, retains cubic Tate sixteen and quartic Tate three, and changes the degree-eight action by exactly the displayed cusp-inner gauge. D.45 proves that this recorded action has no full profinite lift. D.46--D.47 subsequently construct a replacement on their exact relative sources; its full profinite realization remains unresolved. D.47 proves that its full metabelian outer action is outside Galois through the exact cover maps. Its realization by the requested full profinite source groups, and genuine CM-pair attainability in those groups, remain undecided.
\end{proof}

\subsection{The integral signs, CM coordinates, and the genuine congruence action}

\begin{statement}{D.14}{Proved}{What the first integral sign congruences say}
The norm equations give $\epsilon_2=\chi_4(\mu)$ and $\epsilon_3=\chi_{-3}(\mu)$, already at moduli 4 and 3. On the Galois image $\mu=\chi$. Thus these signs add no independent condition to the full universal datum $(\chi,F_g)$: the exact comparison is evaluation of $\chi$ modulo 4 and 3. The unadorned series restricted to $G_K$ has $\chi=1$, and consequently both signs there are one.
\end{statement}
\begin{proof}
For a unit Gaussian norm $a^2+b^2$, exactly one of $a,b$ is odd; its value is one modulo 4. For the Eisenstein norm $c^2+cd+d^2$, reduction modulo 3 is $(c-d)^2$, hence is one for a unit. The signed norm equalities in LM2 C.d2 therefore force the stated signs. Conversely the local norm images are $1+4\mathbb Z_2$ and $1+3\mathbb Z_3$: the norm on sufficiently deep principal units is the trace under logarithm, and inspection at moduli 8 and 9 supplies the first step. These finite images are also enumerated inside \texttt{d\_local\_cm.py}. Evaluating the already specified multiplier proves the asserted factorization; it does not equate an unmarked beta series with its missing multiplier.
\end{proof}

\begin{statement}{D.15}{Known}{The two normalizer coordinates are the projective CM representations}
Use the rigidified modular stack $\mathcal M_{1,1}/\{\pm1\}\simeq[U_{0,3}/S_3]$ and the positive tangential point $\lambda=0$ with tangent1 and the standard branch-cut paths on $0<\lambda<1$. Choose the same etale paths and homology basis in the modular congruence quotient and in LM2 C.d4. Then its coordinates
$T=A^{-1}V$ and $W=\beta^{-1}V$, each modulo its single global sign, are respectively the projective adelic Tate representations of
\[
 E_3:y^2=x^3+1,\qquad E_4:y^2=x^3-x.                       \tag{D.13}
\]
Their determinants are $\chi$; their actions on the CM orders are complex conjugation according to $\chi_{-3}$ and $\chi_4$. Changing the paths conjugates the indicated representations by their transported bases. Quadratic twists of the displayed curves make exactly the same representations modulo the global sign.
\end{statement}
\begin{proof}
The quotient of an elliptic curve by $[-1]$ is a four-pointed projective line, with the identity among the four branch points. Ordering the other three points gives the Legendre parameter and is a torsor under $S_3$. Conversely the double cover of this four-pointed line gives the elliptic curve; its two possible covering automorphisms are exactly $\pm1$. These two constructions are inverse on families and their isomorphisms after rigidification, proving the displayed stack equivalence over $\mathbb Q$.

The projective Tate local system consequently has geometric monodromy $\operatorname{SL}_2(\widehat{\mathbb Z})/\{\pm I\}$. At the elliptic points its inertia generators are the matrices $R,S$ of LM2, satisfying $R^3=S^2=-I$. Transport the Galois action at either elliptic point to the cusp by a fixed path. The resulting path cocycle is a conjugator $A$, respectively $\beta$, with
$u_0(R)=AR^{\chi_{-3}}A^{-1}$ and
$u_0(S)=\epsilon_2\beta S\beta^{-1}$. These are precisely the conjugators in C.d4: their prescribed reductions modulo 2 remove the finite inertia ambiguity. More explicitly, the centralizer of either torsion generator in the free profinite product $C_2\amalg C_3$ is its cyclic factor, whose intersection with the level-two free subgroup is trivial. Hence two such normalized path conjugators coincide. Transport of the Tate action now gives exactly $A^{-1}V$ and $\beta^{-1}V$, proving the formulas.

At $\lambda=1/2$, the change $X=4x-2$, $Y=8y$ gives $Y^2=X^3-4X$, the quadratic twist by 2 of $E_4$. At $\lambda^2-\lambda+1=0$, shifting $x$ by $(1+\lambda)/3$ gives $y^2=x^3+(2\lambda-1)/9$; this constant is the cube of $-(2\lambda-1)/3$, so it is a quadratic twist of $E_3$ over $\mathbb Q(\sqrt{-3})$. The descent to $\mathbb Q$ is fixed by its action on the three nonzero two-torsion points: for $j=0$ the splitting field $\mathbb Q(\sqrt{-3})$ forces the cubic constant to be a rational cube up to a quadratic twist, and for $j=1728$ rational two-torsion forces the coefficient to be a negative square. Thus the descended projective representations are those in (D.13). This is a typed arithmetic comparison, not an inference from the shapes of the two centralizer rings. The CM reciprocity used next is \cite[Corollary 3.4]{dLozano2022}, or equivalently \cite[Theorem 3.10]{dMilne2007}, both opened.
\end{proof}

\begin{statement}{D.16}{Proved}{Exact joint CM image: one additional quadratic relation}
Let $K_3=\mathbb Q(\sqrt{-3})$, $K_4=\mathbb Q(i)$ and $L=K_3K_4$. In the product of the two normalizer groups of D.15 with common determinant and the modulo-two conditions of LM2 C.d4, the joint Galois image has index two. Its exact additional condition is equality of the two restrictions to
\[
 M=\mathbb Q(i,\sqrt[4]{3});\qquad
 K_3^{\rm ab}\cap K_4^{\rm ab}=\mathbb Q^{\rm ab}M.         \tag{D.14}
\]
Thus the two-tori lead has an open image, but the Galois image does not fill the proposed full pair of equal norms. The missing coset is a single $C_2$; the existence of an element of $\widehat{\rm GT}_0$ realizing that coset is a separate unresolved realization question.
\end{statement}
\begin{proof}
Both quadratic rings have class number one, as follows already from their Euclidean norms. CM reciprocity identifies the Weber-function torsion fields with the ray class fields; taking all levels in \cite[Corollary 3.4]{dLozano2022} gives $K_j^{\rm ab}$. The field defined by the Tate representation modulo $\pm1$ over $K_j$ contains these Weber values ($x^3$ for $E_3$, $x^2$ for $E_4$), and is contained in $K_j^{\rm ab}$ since the Tate action over $K_j$ is linear over its commutative CM order. It is therefore exactly $K_j^{\rm ab}$. An element outside $G_{K_j}$ cannot act as $\pm I$, because it conjugates a nonreal CM endomorphism to its distinct complex conjugate. Hence the field of the projective representation over $\mathbb Q$ is also exactly $K_j^{\rm ab}$.

More precisely the image fills the specified individual factor. Reciprocity identifies $\operatorname{Gal}(K_j^{\rm ab}/K_j)$ with $\widehat{\mathcal O}_{K_j}^{\times}/\mathcal O_{K_j}^{\times}$. Reduction modulo 2 maps $\mu_6/\{\pm1\}$ isomorphically to $(\mathcal O_{K_3}/2)^\times$, and $\mu_4/\{\pm1\}$ isomorphically to $(\mathcal O_{K_4}/2)^\times$. Thus removing that finite root of unity gives an isomorphism with the units congruent to one modulo 2, modulo the single global sign. The curves in (D.13) have their full two-torsion over the stated fields, so their projective representations give exactly this isomorphism. The determinant is restriction to the cyclotomic subfield by the Weil pairing. Outside $G_L$, the two conjugation signs are already forced by that determinant and are all realized by $G_{\mathbb Q}$.

It remains to compute the intersection; this step determines the joint image, not merely its openness. Put $E=K_3^{\rm ab}\cap K_4^{\rm ab}$. It contains $\mathbb Q^{\rm ab}$. In $G=\operatorname{Gal}(E/\mathbb Q)$ the two distinct index-two subgroups fixing $K_3,K_4$ are abelian. Their intersection is central, and together they generate $G$. Choosing representatives $x,y$ of their two nontrivial cosets, the closed commutator group is generated by the central element $[x,y]$. Since $x^2$ and $y^2$ are central, this element has order at most two. Therefore $[E:\mathbb Q^{\rm ab}]\leq2$.

The splitting field $M$ of $X^4-3$ has dihedral group of order eight: its automorphisms are $(a,b)$, $0\leq a<4$, $0\leq b<2$, acting by $\sqrt[4]3\mapsto i^a\sqrt[4]3$, $i\mapsto(-1)^b i$, with product $(a,b)(c,d)=(a+(-1)^bc,b+d)$. The subgroups fixing $K_4$ and $K_3$ are respectively $b=0$ and $a+b=0\pmod2$, both abelian. Thus $M\subset E$, and $M$ is nonabelian over $\mathbb Q$. This realizes the bound and proves (D.14). Galois theory now identifies the joint image with the fibre product over this exact intersection. The common-determinant target is the fibre product over $\mathbb Q^{\rm ab}$, so its index is two. The program checks the dihedral group and this finite index explicitly.
\end{proof}

\begin{statement}{D.17}{Proved}{The local indices at two and three}
For $\mathcal O_L=\mathbb Z[i,z]$, $z^2-z+1=0$, the paired local norm map to the equal-norm units of $\mathcal O_{K_3}\otimes\mathbb Z_p$ and $\mathcal O_{K_4}\otimes\mathbb Z_p$ has index two at each $p=2,3$, and is onto at every other prime. Its two local obstruction classes are identified by the global root-of-unity pairs, leaving the single global relation of D.16. Exact finite enumeration gives
\[
\begin{array}{c|rrrrrrrr}
 p^k&2&4&8&16&32&3&9&27\\\hline
 \text{equal-norm pairs}&6&48&384&3072&24576&24&648&17496\\
 \text{norm image}&6&24&192&1536&12288&12&324&8748.
\end{array}                                               \tag{D.15}
\]
\end{statement}
\begin{proof}
At either bad prime one quadratic subfield is unramified and the other ramified. Given a compatible pair $(a,b)$, first lift its coordinate in the unramified subfield through the ramified quadratic norm. It is a norm because the quadratic Hilbert-symbol obstruction equals the Hilbert symbol of its rational norm, and that rational norm is already a norm from the ramified quadratic subfield. Its valuation is zero, so a lift is a unit. The remaining discrepancy has norm one in the ramified subfield. Hilbert--90 writes it as $c/\bar c$. It can be corrected by a unit from the biquadratic field exactly when the valuation of $c$ in that ramified subfield is even: norms from its unramified quadratic extension have even valuation, and multiplying $c$ by a rational scalar changes its valuation by an even integer. This gives precisely one $C_2$ obstruction. The unit norm in an unramified extension is onto, so no further obstruction remains.

The pairs $(-1,1)$ and $(1,i)$ represent the nonzero class at each prime. The identity
$N_{L/K_3}(z+i\bar z)=-1$, $N_{L/K_4}(z+i\bar z)=-i$
shows that their sum is zero. At other primes the integral norm torus has connected smooth kernel; norm surjectivity on residue fields and on successive principal-unit quotients proves surjectivity. The code enumerates the rings directly: for $a+bz+i(c+dz)$ its two norms are
$(a+bz)^2+(c+dz)^2$ and
$(a+ci)^2+(a+ci)(b+di)+(b+di)^2$.
It tests invertibility by the rational norm and records the entire image size. The signs above are also tested at every displayed level. The global roots of unity contribute the same nonzero class at both bad primes, hence remove only the diagonal $C_2$. This agrees with, and does not replace, the field-intersection proof of D.16.
\end{proof}

\begin{statement}{D.18}{Known}{The genuine cusp congruence action retains the full rank-one Kummer coordinate}
In the positive tangential normalization $\lambda=0$, tangent1, take $q^{1/n}$ with positive real leading coefficient $256^{-1/n}$ and $\zeta_n=\exp(2\pi i/n)$. Use the resulting Tate basis, matched to the Legendre cycles over $[0,\lambda]$ and $[\lambda,1]$, with positive local monodromy $\left(\begin{smallmatrix}1&2\\0&1\end{smallmatrix}\right)$. The genuine normalized congruence representation is
\[
 V_g=\begin{pmatrix}\chi(g)&-8\kappa_2(g)\\0&1\end{pmatrix}.            \tag{D.16}
\]
Compatible roots may be chosen so that
$\kappa_{-1/2}=(\chi-1)/2-\kappa_2$. Consequently the continuous inverse comparison on this image is
\[
 (\chi,V_{12})\longmapsto
 \left(\chi,\frac{\chi-1}{2}+\frac{V_{12}}8\right).          \tag{D.17}
\]
Multiplication by eight is a homeomorphism from $\widehat{\mathbb Z}$ onto $8\widehat{\mathbb Z}$, so no part of this Kummer coordinate is lost in the genuine normalized congruence action. In the basis of the manuscript with local monodromy $\left(\begin{smallmatrix}1&-2\\0&1\end{smallmatrix}\right)$, conjugation by $\operatorname{diag}(-1,1)$ reverses the off-diagonal sign. On \emph{outer} congruence actions this translation is forgotten: the comparison retains only $\chi$ modulo squares.
\end{statement}
\begin{proof}
The Legendre invariant is
$j(\lambda)=256(1-\lambda+\lambda^2)^3/(\lambda^2(1-\lambda)^2)$.
Inverting $j(q)=q^{-1}+744+\cdots$ gives
$q=\lambda^2/256\,(1+\lambda\mathbb Q[[\lambda]])$.
Tate uniformization gives the exact torsion sequence with basis $\zeta_n,q^{1/n}$ and matrix $\left(\begin{smallmatrix}\chi&\kappa_q\\0&1\end{smallmatrix}\right)$; see \cite[\S11, Theorem 11.17 and the preceding uniformization]{dJoshiTate}. For the formal field used here, the comparison follows by base change of the universal Tate equation and its torsion representatives over $\mathbb Z((q))$: the convergent uniformization identities are identities of Laurent series with integral coefficients, and remain identities after $q=q(\lambda)$. Thus the torsion representatives are $\zeta_n^a q^{b/n}$ over $\overline{\mathbb Q((\lambda))}$. The tangential section fixes compatible $\lambda^{1/n}$ and the unique roots with constant coefficient one of the last power-series unit. Its only remaining Kummer action on $q^{1/n}$ is therefore that on $256^{-1/n}$, namely $-8\kappa_2$. The Legendre curve is a quadratic twist of the split Tate curve at this cusp; that twist multiplies the entire matrix by a single sign and disappears in the projective congruence local system. Normalizing its lower-right entry to one yields (D.16). To fix the residual shear, let $K(m)=\int_0^{\pi/2}(1-m\sin^2\theta)^{-1/2}d\theta$. The branch-cut periods have ratio $\tau=iK(1-\lambda)/K(\lambda)$, purely imaginary and positive for $0<\lambda<1$. Hence $q=\exp(2\pi i\tau)$ and its chosen roots are positive real. A shear would replace $q^{1/n}$ by $\zeta_n^b q^{1/n}$ and change the off-diagonal entry by $b(1-\chi)$; positivity for every $n$ forces $b=0$. This proves that the Tate basis is the specified global Legendre branch-cut basis, rather than merely a basis with the same local monodromy.

Choose $(-1)^{1/n}=\zeta_{2n}$. Its Kummer cocycle is $(\chi-1)/2$, which proves (D.17) by multiplication of root systems. The stated change of basis explicitly handles the orientation convention. Finally conjugation by $V_g$ on the special linear congruence quotient has outer class determined by its determinant modulo squares, as established in Part A; its unipotent upper-triangular factor is inner. Thus the often proposed claim that the rank-one translation is visible only metabelianly is false for genuine normalized actions, while the precise outer quotient does discard it.
\end{proof}

\begin{statement}{D.19}{Open}{The residual GT and full profinite realization questions}
D.1--D.4 settle item3, including general openness, the quadratic failure of surjectivity, joint independence, and the local inertia images. D.15--D.18 settle the stated Galois comparison of the two CM coordinates and the rank-one congruence coordinate. The unresolved parts of D(e) and D-after-LM2(1) concern the full profinite GT beta/gamma presentation and actual realization outside Galois: whether the extra leading metabelian directions of D.11 lift from $\widehat{\rm GT}_0$ or the full profinite loop-preserving group, and whether $\widehat{\rm GT}_0$ supplies the missing index-two CM coset of D.16. These assertions remain on the coverage ledger; none follows from solving only matrix or Lie equations.

D.39 extends a new representative through the entire all-prime mixed Fermat homology tower, jointly with the retained degree-eight quotient. Its common quotient stabilizes at exponent eight. The exact two-adic system has cokernel killed by sixteen, and the half-average gives every odd-prime component. The resulting multiplier-one automorphism retains the normalized G6 and T3 actions, replaces the prior J3 and J12 actions by the specified new components, and retains its checked J48 seed at every later level. Both compact CM matrices and the fixed-fibre pair survive. The exact source kernel is the intersection of the closed L-commutator subgroup with D9K; the nonzero degree-sixteen word remains in it. The projection to the earlier finite joint group is strict, detected by $x^{144}$. D.45 proves that this recorded action has no full profinite lift. D.46--D.47 subsequently construct a replacement on their exact relative sources; its full profinite realization remains unresolved.

D.40 extends that action through the actual image in every finite
perfect quotient of the marked source. Their exact common quotient
with the retained tower is $C_6$, and the normalized conjugators pair
with the identity on all perfect coordinates. A further actual
$S_7$ quotient gives the fibre product over the common sign $C_2$.
The exact remaining full-source kernel maps onto $A_{10}$ under the
computed $S_{10}$ quotient. These maps retain both compact-cover and
fixed-fibre CM coordinates; D.45 excludes a full profinite lift of this recorded action.

D.41 then extends the action through all actual symmetric quotients
simultaneously. Their even image is the full product of the
alternating factors indexed by distinct marked kernels, and the
common quotient with D.40 is exactly the sign $C_2$.
Both projection kernels are computed. The remaining full-source
kernel maps onto $\operatorname{PSL}_2(\mathbb F_7)$ under the
specified $\operatorname{PGL}_2(\mathbb F_7)$ map and kills the
canonical perfect mod-seven coordinate. This retains the actual
CM data; D.45 excludes a full profinite lift of the recorded action.

D.42 extends this actual source action through the full relative
pronilpotent kernel of the specified projective-linear quotient.
Its maximal common marked quotient has base order $54$ and kernel
$\mathbb Z_3^{10}$. The nonzero common conjugator is retained by
an actual coprime centralizer lift, and every projection kernel is
computed. The remaining source kernel maps onto
$\operatorname{PSL}_2(\mathbb F_{11})$ under the specified
$\operatorname{PGL}_2(\mathbb F_{11})$ map. D.45 supplies a finite sphere quotient that excludes every full profinite lift of this recorded action.


D.43 extends the retained action simultaneously through all relative
pronilpotent completions whose marked finite base has nonabelian
simple even subgroup. The exact common quotient with the retained
tower is $C_{54}\times_{C_6}D_{\mathscr J}$, with its actual
source-kernel formula. One compatible centralizer lift supplies the
conjugator over the whole family. The exact remaining source kernel
maps onto $A_7\times A_7$ under an explicit quotient whose involution
exchanges the two simple factors. D.45 excludes every full profinite lift of this recorded action.

D.44 extends the same quotient action to the complete family of
anabelian even bases with their relative pronilpotent kernels.
The common quotient and both source kernels are exact, and an
actual $A_5^4\rtimes S_4$ marking detects residual kernel $A_5^4$.
D.45 then excludes a genuine full profinite lift of the entire
recorded normalized D.39 action and every D.40--D.44 extension
retaining it. Its sphere projection through $B/D_9F$ has a
nonzero degree-nine obstruction after every lower centralizer
and multiplier freedom is included. The elliptic projection
has an explicit corrected lift to $B/D_{10}B'$. D.46--D.47 subsequently supply a replacement realizing the prescribed
CM pair on their exact relative sources. The full source image
questions in D.13 and D.19 remain to be determined.

D.46 constructs a replacement on the entire relative pronilpotent
completion of the actual rank-97 subgroup $N$. Its full relative
pro-two component is an actual arithmetic commutator action with
identity sphere projection and compact elliptic scalar $\eta_*$;
every odd relative-prime component is the identity. The two actions
join through explicit coprime centralizer lifts. The resulting
$G_2,G_3$ actions determine the prescribed pair $T_4=I,W_3=[i]$,
and the reconstructed parameter has an actual $\widehat F'$ preimage.
Thus the replacement solves the entire rank-four relative pro-two
lifting step. Its exact remaining full-source kernel has image
$A_7$ under the displayed $S_7$ quotient. Its full profinite
realization remains in the lifting problem.

D.47 extends this replacement coherently over every Fermat level.
For \(M=\widehat N\cap\widehat F'\), it proves both inclusions
identifying the closed relative pro-prime source with the inverse
system over the actual \(L_m\). The common marked quotient is
\(\mathcal H_\infty\), and the simultaneous kernel is
\(\prod_pM^{(p)}\). The source kernel lies in \(M'\subset
\widehat F''\), so the action descends to the full metabelian
quotient, with every profinite deck coordinate retained. Its
coefficient-two action is the actual arithmetic \(h\) action;
every odd coefficient action is identity. The actual marked
cubic-four and quartic-three cover maps give \((-I,I)\), outside
the proved Galois joint graph. The Heisenberg quotient and full
Alexander-module argument also proves this in the outer action
space. All comparison kernels to D.46 are computed. The remaining
full-source kernel still has exact \(A_7\) image in the displayed
\(S_7\) quotient. The full GT0 and full loop-preserving
\(\widehat F\) lifting questions remain unresolved.

D.48 extends the entire D.47 action over the central double cover
\(\mathcal E=\widehat B/[\widehat N,\widehat B]\), jointly with
every Fermat level. The new subgroup \(M^+\) has index two in
\(M\); its relative pro-two source is unchanged. The comparison
kernel is the product of the explicitly identified odd-prime
kernels. The actual \(h\) action fixes the enlarged base, and
the coprime centralizer corrections retain the complete D.47
action and parameter. On the new finite mixed quotient
\(\mathcal J_3\), the exact old source-kernel image is
\(C_3^{96}\); the new source kernel has trivial image there.
Its full common quotient with D.47 is
\(\mathcal E\times_QG_3\), and the identity action on
\(\mathcal J_3\) agrees with the retained action on that whole
quotient. The remaining full-source kernel has exact
\(A_7\) image; full profinite realization remains unresolved.


D.49 extends this same action over the maximal marked quotient
fixed pointwise by the actual arithmetic \(h\), and over every
relative-prime source above that base. Its exact comparison
with D.48 retains both kernel layers, including an actual
\(A_5\) quotient of the base residual. The full congruence
quotient \(\operatorname{SL}_2(\widehat{\mathbb Z})/\{\pm I\}\),
with one global sign, factors through the identity-three
component; the new action is therefore pointwise identity
on this entire quotient. Its remaining source kernel has
exact image \([V_{3,2},A]\) in the specified mixed homology
quotient over \(G_3\). It contains an actual cyclic moving
subspace of dimension \(64\cdot3^{97}\); the dimension
of the full arithmetic span \(A\) is retained as unevaluated.
The earlier central iteration stabilizes by the exact
integral equality \([L,\widehat B]=L\).
The remaining mixed source and full profinite realization
questions stay active.


D.50 extends this same retained action through the entire
two-adic homology tower over the actual marked \(G_3\).
The exact residual image is the integral moving summand
\(U^-\), and every common quotient is computed from the
original source kernels. A regular-edge Fitting projection
and its explicit crossed homomorphism repair the initial
peripheral endomorphism to an automorphism; the cyclic
norm correction fixes the original cusp. Compatible
conjugators and invertible reduction modulo two give all
higher lifts and their continuous inverse limit. The
full parameter, CM pair, metabelian and congruence actions
are retained. The remaining source kernel has a nonzero
image containing \(W_1\wedge V\) in the next actual
class-two exponent-four quotient. Extending across that
remaining kernel and attaining the full profinite source
remain active steps.


D.51 retains this entire chosen homology action and both
full conjugators and extends them through the moving
part of the next actual class-two exponent-four quotient.
The exact old source image is \(Z=W\wedge V\); its
invariant and moving dimensions are respectively
\(m^2(q-1)/2\) and \(w\rho/2\). The actual common
quotient is \(\mathcal T_2/Z\). Characteristic descent,
exterior preservation of \(Z^A\), and Frattini generation
prove the new action on the joint source. Its projection
adds the whole moving central summand and retains every
previous parameter, CM, metabelian and congruence comparison.
The remaining exact image is \(Z^A\ne0\). In the separate
order-1080 verification, the complete corrected cusp word
has 4,159,444 letters. That example's fixed homology matrix
has intersection-pairing defect of rank 72, excluding its
extension to the unquotiented class-two peripheral source.
The arithmetic invariant obstruction and the full
profinite realization remain to be evaluated.


D.52 solves the entire invariant boundary equation in the
separate marked order-1080 source. Its 120-dimensional
space consists of allowable conjugator pairs. A retained
weight-60 parameter solves all thirteen actual quadratic
equations; the complete 6,352,056-letter corrected word,
both torsion relations and generation are verified.
The new homology action has deviation rank 96 over
\(\mathbb F_2\) and preserves the mod-two ribbon form.
It replaces the excluded finite representative while
retaining the actual invariant quotient. It also has a
complete two-adic homology continuation with the entire
new class-two action retained, exact common quotient and
both projection kernels. Its next exponent-eight residual
is twice the commutator module, witnessed by a nonzero
272-letter commutator square. The arithmetic \(G_3\)
correction and the full profinite realization remain to
be evaluated through their own actual source maps.


D.53 computes the full actual moving peripheral-pair space,
its exact homology image and kernel, and a positive arithmetic
dimension bound. Its integral lifts give actual doubled
automorphisms of the complete D.50 source. D.54 lifts every
such difference to the full joint class-two source, through
an explicit compatible square module. The fixed arithmetic
cusp defect is preserved literally, and the actual affine
quadratic map descends through every same-action ambiguity.
Its arithmetic zero locus and full profinite realization
remain to be evaluated. In the separate order-1080 source,
the full finite homology zero locus has exactly 864 elements;
all 262,144 parameters, the original equations, generation
and the complete long cusp words have been checked.


D.55 gives the exact full-group norm from the arithmetic
cusp obstruction to the compact orientation-bivector defect.
Its local image lies in the norm kernel, and its residual
quotient has explicit Tate-norm and compact-projection maps.
Every point of the reduced boundary zero set has invertible
homology, so generation follows there. The actual ordinary
relative-$H_2$ orientation comparison is an isomorphism.
The arithmetic zero set and residual norm kernel remain
unevaluated, with all maps retained for the next attack.


D.56 evaluates an actual irreducible $C_3^4$ quotient of
the arithmetic normal span and the unique parameter $(2,1,2,1)$.
The marked quotient has order 3888 and actual action
$\operatorname{Int}(t^6)$. Its exact $A_4$ comparison kernel
gives surjective integral homology and class-two boundary maps.
The smaller arithmetic correction space has dimension 40,
with invariant central target dimension 2560. Its actual
quadratic zero set and full profinite realization remain
unevaluated; every source and comparison kernel is retained.


D.57 now solves the complete ten-coordinate boundary equation
on the actual order-$3888$ arithmetic quotient, for both integral
rank-nine branches and all sixteen framed cases. The literal
genus-three cover and degree-eight elliptic map, exhaustive integral
conjugator search, separate plus-central alignment and augmentation
correction give the asserted automorphism of
$\widehat B/(R_h\cap M_*\cap L_{*,2})$.
Its entire $W_h$ action and prescribed normalized parameter are
retained. The exact remaining quotient kernels carry the full
$G_3$ invariant and full profinite realization questions onward.


D.58 retains the entire D.57 action and its conjugators while
restoring the full integral arithmetic homology. The actual
projective repair stays in
$\mathcal L=\ker(U^-\to U_*^-)$.
It then extends through the remaining moving class-two module,
with exact residual kernel
$\mathcal K_0=\ker(Z\to Z_*)^A$.
All $2560$ solved invariant coordinates are retained; the remaining
dimension is $m^2(|A|-1)/2-2560$.
The next actual invariant boundary and the full profinite realization
remain to be evaluated across this specified source map.


D.59 corrects the full homology on every active cusp-safe free
character orbit of the actual arithmetic source. Its invariant
central block has an injective full norm and compactification;
the exact unquotiented orientation identity therefore kills its
boundary. The residual kernel has dimension
$48w-2560-5\,308\,416\kappa$, where
$\kappa\geq830\,900\,448\,697\,714\,719\,719$.
It contains the specified $366\,080$-dimensional small-character
submodule. Its remaining boundary and full profinite realization
are the next actual-source steps.

D.60 now repairs twelve further central blocks on the same
arithmetic source. Its integral projective-cover construction
retains the complete D.57 action and the prescribed $\mathcal W_h$
projection, while changing the D.59 homology action on two
explicit character orbits. The full norm and compactification
are injective on the new $12\,288$-dimensional central block,
so its twenty-eight boundary coordinates vanish by the
unquotiented orientation identity. The explicit residual
small-character module has dimension $353\,792$ and
$612$ coinvariants; its actual remaining boundary and the
full profinite realization remain to be evaluated.

D.61 makes the exact integral homology replacement on the two
size-twelve character orbits. It retains D.57 and $\mathcal W_h$,
and its injective norm and compactification kill the actual
orientation boundary on seven previous self blocks and twelve
new cross blocks. The resulting lift retains $22\,016$ further
central coordinates and fifty coinvariants. Its specified
remaining module has dimension $331\,776$ and $562$ coinvariants.
At the D.61 stage the four actual involution-cover divided norms
leave twenty-six named self-wedge classes. D.62 evaluates their
arithmetic correction by the actual source words below.

D.62 cancels all twenty-six projective self-wedge classes by
actual source-involution and infinite-cusp correcting words.
The full cusp wrap is retained; its extra commutator is killed
by the specified projector. The literal comparison on $Y$ is
$\operatorname{Int}(b_{t,Y})\alpha_{61,Y}$, with unchanged
homology $T_{61}$ and retained D.57 and $\mathcal W_h$ maps.
The new lift retains $6656$ further central coordinates and
twenty-eight coinvariants. Its specified residual has dimension
$325\,120$ and $534$ coinvariants. The arithmetic boundary on
this residual and the full profinite source realization remain
the active original questions.

D.63 orthogonalizes the actual compact projective summands
through its explicit integral graph and lifts them inside the
retained source kernel. The actual peripheral derivation gives
$T_{63}=(1-e)T_{61}(1-e)+e$ and the exact lower action
$\operatorname{Int}(b_{t,Y})\beta_Y$, retaining the complete
D.57 and $\mathcal W_h$ maps and all old projector kernels.
The new lift adds $40\,960$ central coordinates and fifty-two
coinvariants. Its specified remaining kernel $\mathcal J_{63}$
has dimension $284\,160$ and $482$ coinvariants; its compact
norm has rank six and kernel dimension $476$.
The arithmetic boundary on this literal residual and the
full profinite source realization remain the original active
questions.

The successive residual descriptions above refer to their
numbered stages. D.64 extends the actual action through the
onto quotient $\Phi_{64}$ of $\mathcal J_{63}$ of dimension
$1930$, with fourteen coinvariants. Marked compact isometries
and actual source-involution words give
$T_{64}=T_{63}+A_1m_1+A_4m_4$ and its literal lower-action
comparison, retaining both D.57 branches, all sixteen framings,
the full $\mathcal W_h$ maps and source conjugators.
Its post-word moving correction uses the actual
$R_{50}\cap L_{*,2}$ source lifts before the final augmentation.
The exact remaining kernel $\mathcal J_{64}=\ker\Phi_{64}$
has specified dimension $282\,230$.
The fourteen uncovered compact-test coordinates retain
their test scope; the entire residual arithmetic boundary
and the full profinite source realization remain active.

D.65 retains the full second and third compact blocks by an actual
integral source correction. Its signed exterior norm divided by two
has rank six in each ten-dimensional test. Coinduction realizes the
full image, and the actual source-involution words kill both rank-four
kernels. The further quotient has dimension $5752$ and twelve
coinvariants, with total retained dimension $7682$ and twenty-six
coinvariants. Both D.57 branches, literal small conjugators and the
full $\mathcal W_h$ maps are retained. The lower action is exactly
$\operatorname{Int}(b_{t,Y})\beta_Y$ with its stated derivation and
all thirty-six cusp factors. The remaining specified kernel has
dimension $276\,478$ and maps onto the old compact norm kernel of
dimension $1918$. Its two remaining tested coordinates retain that
test scope. The residual source boundary and full profinite
realization continue to require their actual lifts.

D.66 fills the complete fourth compact comparison block by actual
quadratic source automorphisms. Its integral projective cover and full
quadratic source square realize both remaining compact directions,
retaining the literal small conjugators, both D.57 branches, all
$\mathcal W_h$ maps and the preceding three compact blocks. The total
retained compact quotient has dimension $9600$ and $28$ coinvariants.
The specified remaining source kernel is $J_{66}$, of dimension
$274\,560$. Its cusp-radical correction and the full profinite
realization continue to require their actual source maps.
\end{statement}
\begin{proof}[Concrete residual instances]
The first extra metabelian special vector occurs at weight five, where the free special source has dimension three and its invariant part dimension one; exact preimages are in \texttt{d\_metabelian.json}. The first extra invariant metabelian vector occurs at weight nine. D.12 supplies their pro-$\ell$ realizations and identifies the nonnilpotent-quotient obstruction to transferring that construction. For the CM problem, D.17 supplies an explicit finite matrix coset, represented locally by $(-1,1)$ at one bad prime and the identity at the other, whose two obstruction classes disagree. D.16 proves that this coset has no Galois lift. Its possible lift to an actual full profinite $\widehat{\rm GT}_0$ element must solve the conjugator and generation equations before congruence reduction; the matrix reconstruction of LM2 C.d4 does not do so. The source search reached the original Ihara Part II through the RIMS1224 PostScript file, and the exact modern profinite lift criterion through Dolgushev--Guynee. Part I's RIMS1177 catalogue entry has no download; the PDF, PostScript and DVI locations returned 404. D.20--D.22 record the additional results actually obtained. D.25--D.33 carry out finite and infinite quotient tests for the missing CM coset: the twenty-three initially selected nonidentity shadows and two later central representatives are excluded, while successive explicitly distinguished replacements supply actual compatible actions on the stated quotient systems. D.31 constructs an actual action on $\mathcal S_2^{(7)}$, combined with all odd relative pro-prime kernels and D.28's mixed quotient through its exact common quotient. D.32 proves that its recorded $T_2$ action is obstructed at the next degree, including every same-action conjugator and multiplier choice. D.33 then changes the lower action explicitly and supplies a compatible degree-eight replacement with a class-four kernel, including the whole class-two tower and the same mixed-prime quotient. The remaining relative-source kernel is exactly $\gamma_3N^{(2)}\cap D_9(K^{(2)})$, and an actual nonzero triple commutator is exhibited in it. The entire CM pair has not been excluded. D.34--D.35 give exact structural maps for the relative problem: the marked elliptic cover, the full outer deck centralizer, actual odd-deck boundary lifts, and a genuine arithmetic surjection onto the single compact CM torus. The inverse image of the recorded D.33 action under (D.35.5a) retains the full finite coordinates modulo the specified cusp-inner subgroup. D.45 proves this particular fibre empty, using the full torsion and quotient-preservation maps and an actual cusp-inner transport. D.36 gives the actual quartic-cover and common cyclic-twelve maps, the arithmetic joint degree-six fibre, and the retained gauge-invariant sign discrepancy. The opened cover-to-fixed-fibre source equations retain their stated Galois scope. D.37 supplies a further exact mixed Fermat-twelve shadow, retaining D.33 on $G_6$ and evaluating both actual cover matrices as $(-I,I)$. Its universal relative-prime quotient maps and mixed kernel are explicit. D.38 then proves compatibility for a new Fermat action with the retained degree-eight component, through their exact maximal common quotient. The simultaneous finite group, both projection kernels and a nonzero word in the remaining full-source kernel are computed. The unresolved steps are the complete lift tests for the two indicated metabelian vectors and for the missing CM coset.
\end{proof}

\subsection{The original profinite beta/gamma relations and their graded consequence}

\begin{statement}{D.20}{Known}{Ihara's exact cyclic-cover compatibility and gamma multiplication}
Write a full profinite GT element as $(\lambda,f)$ acting by
$x\mapsto x^\lambda$, $y\mapsto f^{-1}y^\lambda f$.
For $n\geq1$ let $H_n=\ker(\widehat F_2\to\mathbb Z/n)$, where $x\mapsto1$, $y\mapsto0$, and define the continuous epimorphism
\[
 \phi_n:H_n\longrightarrow\widehat F_2,\qquad
 x^n\mapsto x,\quad y\mapsto y,\quad x^{-j}yx^j\mapsto1\quad(0<j<n).
\]
For every $g\in G_{\mathbb Q}$ one has
\[
 \phi_n(f_g)=y^{-\kappa_n(g)}f_g,\qquad
 \phi_n\circ g|_{H_n}=g\circ\phi_n .                         \tag{D.18}
\]
Here the roots of the positive integer $n$ use the positive real branches of D.2. These are \cite[Theorem 1 and Proposition 1]{dIhara2000}.

For every $\sigma=(\lambda,f)\in\widehat{\rm GT}$, put
$c_n(f)=\operatorname{exp}_y(\phi_n(f))\in\widehat{\mathbb Z}$.
The continuous cover maps above satisfy
\[
 \phi_n(f)=y^{c_n(f)}f,\qquad
 \phi_n\sigma|_{H_n}=\sigma\phi_n,\qquad
 c_{mn}(f)=c_m(f)+c_n(f).
 \tag{D20.1}
\]
The full profinite assertion follows from
\cite[Theorem8.1 and Remark8.6]{dEnriquez2007}; the exact
comparison of its generators with ours is proved below.
On the Galois subgroup $c_n(f_g)=-\kappa_n(g)$.

The full profinite beta function of $\sigma\in\widehat{\rm GT}$ has Ihara's factorization
\[
 B_\sigma(s,t)=\frac{\Gamma_\sigma(s)\Gamma_\sigma(t)}{\Gamma_\sigma(s+t)},
 \qquad s,t\in\mathbb Q/\mathbb Z.                          \tag{D.19}
\]
The coefficient ring and completed group-ring interpretation are those of the hyperadelic gamma function; equality precedes every prime-power specialization. For Galois elements Anderson's multiplication formula, in Ihara's corrected convention, is
\[
 \frac{\prod_{nc=0}\Gamma_g(s+c)/\Gamma_g(c)}{\Gamma_g(ns)}
       =1\otimes\exp(-2\pi i\,ns\kappa_n(g)).                 \tag{D.20}
\]
Ihara Part II, Introduction and \S2, states (D.19) and (D.20), attributes (D.20) to Anderson's Corollary 8.6.3, and gives the following exact relations among its intermediate subgroups:
$G_{\mathbb Q}\subset\widehat{\rm GT}_K\subset\widehat{\rm GT}_{A'}$ and
$G_{\mathbb Q}\subset\widehat{\rm GT}_A\subset\widehat{\rm GT}_{A'}$.
Here $K$ means all the commuting squares in (D.18), $A$ means (D.20), and $A'$ means the two distribution identities in the proof below.
Combining these proved inclusions with (D20.1) gives the
stronger equality
\[
 \widehat{\rm GT}_K=\widehat{\rm GT}_{A'}
                   =\widehat{\rm GT}.
 \tag{D20.2}
\]
\end{statement}
\begin{proof}
Ihara's earlier original notes \cite[Sections2--4, pp.46--52]{dIhara1997Kokyuroku} give the profinite beta definition, the associativity statement yielding gamma factorization, and a determinant proof sketch for multiplier one. We read all ten pages. The general normalization here retains the opened Part II and Anderson comparisons.

The original RIMS1224 version was downloaded, decompressed, converted from PostScript, and read. Its proof of Theorem 1 identifies $H_n$ with the fundamental group of $\mathbb P^1-\{0,\mu_n,\infty\}$, and $\phi_n$ with the map induced by filling the punctures different from 1, after the change $t\mapsto t^n$. Coefficientwise Galois action on Puiseux series commutes with this change, proving the second equality of (D.18). Applying the two sides to $y$ shows that $\phi_n(f_g)f_g^{-1}$ centralizes $y^{\chi(g)}$, hence is a power of $y$. Its exponent is the logarithmic Kummer quasi-cocycle $\Psi_n^{(0)}(g)$, characterized in the original Theorem 1 by
$(n^{-r})^{g-1}=\exp(2\pi i r\Psi_n^{(0)}(g))$ for $r\in\mathbb Q$.
Thus $\Psi_n^{(0)}=-\kappa_n$, proving the first equality with its sign.

We now give the full-GT reduction on these same profinite groups.
Enriquez uses the Schreier generators
$x^n,x^j yx^{-j}$ and a filling map $\delta_n$ killing
the generators with $0<j<n$. The identity
\[
 x^{-j}yx^j=x^{-n}(x^{n-j}yx^{-(n-j)})x^n
\]
shows that its continuous extension is exactly $\phi_n$.
The cited theorem therefore gives the first equality in
(D20.1), with precisely our $c_n$.
Its proof uses the actual cyclotomic configuration covers,
the commuting strand-doubling maps of Lemmas8.2--8.4,
and the pentagon. Lemma8.5 removes the central factors;
erasing strand zero gives $\delta_n(f)=y^{c_n(f)}f$.

Here is the resulting intertwining map explicitly.
Since $\lambda$ is a unit modulo $n$, $\sigma$ preserves
$H_n$. On $x^n$, both maps $H_n\to\widehat F_2$ give
$x^\lambda$; on $y$, both give $f^{-1}y^\lambda f$.
For $z_j=x^{-j}yx^j$, write
$j\lambda=a+nb$ with $1\leq a<n$ and
$b\in\widehat{\mathbb Z}$. Then
\[
 \sigma(z_j)=k^{-1}x^{-nb}z_a^\lambda x^{nb}k,
 \qquad k=x^{-j\lambda}fx^{j\lambda}\in H_n.
\]
Applying $\phi_n$ gives $1$, also equal to
$\sigma\phi_n(z_j)$. These free profinite generators
prove the second equality.

Finally $\phi_n(H_{mn})\subset H_m$ and
\[
 \phi_m\phi_n|_{H_{mn}}=\phi_{mn}:
                       H_{mn}\longrightarrow\widehat F_2.
\]
Indeed $x^{mn}$ maps to $x$, while each
$x^{-j}yx^j$ with $0<j<mn$ is killed at the first step
when $n\nmid j$, and at the second when $j=nk$,
$0<k<m$. Applying this identity to $f$, using
$f\in\widehat F_2'\subset H_{mn}$, gives
$y^{c_{mn}(f)}f=y^{c_n(f)+c_m(f)}f$.
The embedded procyclic subgroup generated by $y$ has
exponent group $\widehat{\mathbb Z}$, proving additivity.
Ihara's Proposition2 now applies to every full-GT element,
and the inclusions already displayed give (D20.2).

For completeness the identities defining $A'$ are
\[
 \Psi_{mn}^{(0)}=\Psi_m^{(0)}+\Psi_n^{(0)},\qquad
 \Psi_{mn}^{(in)}=\Psi_m^{(i)}\quad(i\ne0\bmod m).
\]
The functions $\Psi_n^{(i)}$ are defined by the action of $f$ on
$(1-\zeta_n^{-i}t^{1/n})^r$; Proposition 2 proves these identities directly from (D.18). Proposition 3 proves that they are exactly the logarithmic derivative of (D.20), using Part I's Theorem 2, whose formula is reproduced there. The original Introduction explicitly says that the proof of (D.19) uses the five-cycle relation. These statements specify which equation is used and which group it constrains.

Anderson's original Corollary 8.6.3, p.103, was also visually
checked. With $\mathcal W=\prod_pW(\overline{\mathbb F}_p)$
from his \S7.1, it states, in
$(\overline{\mathbb Q}\otimes\mathcal W)^\times$,
\[
 \prod_{nb=a}\Gamma_g(b)=
 \bigl((n^{\langle-a\rangle})^{g-1}\otimes1\bigr)
 \Gamma_g(a)\prod_{nc=0}\Gamma_g(c).
\]
The roots of $n$ are positive real. Put $a=ns$ and divide
by the two displayed products; their Kummer factor is
$\exp(-2\pi i\,ns\kappa_n(g))\otimes1$.
Ihara's coefficient convention places $\mathcal W$ first;
the tensor flip $z\otimes w\mapsto w\otimes z$ gives exactly
(D.20), including its sign.

There is a precise obstruction to reversing the logarithmic derivative in the hyperadelic ring. For complementary nonzero idempotents $e,e'\in\widehat{\mathbb Z}$, $ee'=0$, the nonconstant element $e x^{e'}$ has derivative zero. It is nonconstant already after projection to a coefficient prime supported on $e$ and a finite deck quotient supported on $e'$. Thus the kernel of the derivation contains more than constants. Ihara's final Remark explicitly records that $\widehat{\rm GT}_K\subset\widehat{\rm GT}_A$ was not proved there. Formula (D20.1) establishes the full-GT cover and distribution identities; the gamma multiplication formula (D.20) retains its proved Galois scope.
\end{proof}

\begin{statement}{D.21}{Proved}{Galois and full GT have the same rational metabelian graded image through weight ten}
For every prime $\ell$, the rational graded images of the actual groups $G_{\mathbb Q}$ and $\widehat{\rm GT}$ in the metabelian outer action, in weights 1 through 10 and after removing the cyclotomic multiplier, both have dimensions
\[
                 0,0,1,0,1,0,1,0,1,0.                     \tag{D.21}
\]
In weight $n\geq3$ odd the common line, expressed in the $p(X,Y)[x,y]$ coordinate of D.11, is
\[
 \mathbb Q_\ell\,
 \frac{(X+Y)^n-X^n-Y^n}{XY}.                                \tag{D.22}
\]
The inclusion of this line in the $\widehat{\rm GT}_0$ upper-bound space is the explicit polynomial inclusion in D.11. The latter upper bound has an additional line in weight nine.
\end{statement}
\begin{proof}
We give the graded comparison explicitly. Map the free profinite group to its pro-$\ell$ quotient, take logarithms in each rational nilpotent quotient, and then take the first nonzero homogeneous term of a filtered GT element. Expanding the two-cycle, three-cycle and five-cycle equations gives their infinitesimal equations for this term. The infinitesimal pure-braid generators are the initial classes of the logarithms of the pure-braid generators. Their relations are
$[t_{ij},t_{kl}]=0$ for four distinct indices and
$[t_{ij},t_{ik}+t_{jk}]=0$; the successive split fibrations forgetting a strand identify their free Lie kernels, so these relations give the associated graded pure-braid Lie algebra used in the expansion. Thus the displayed map sends the leading term into $\mathfrak{grt}_1(\mathbb Q_\ell)$.

Furusho's \cite[Corollary 0.3, Theorem 0.5, Proposition 5.1 and Lemma 5.2]{dFurusho2011} gives the gamma constraint on its metabelian image. Here is the comparison with our coordinate, including the factor that matters. For a Lie element $h=h_xx+h_yy$, put $M(h)=(h_yy)^{\rm ab}$. Direct expansion gives
\[
 M([x,y])=XY,\quad M([x,h])=XM(h),\quad M([y,h])=YM(h)
 \quad(h\in L').
\]
It follows that $M(p(X,Y)[x,y])=XYp(X,Y)$ and that $M$ kills $L''$. Furusho's gamma constraint is
$M(h)=\gamma(X)+\gamma(Y)-\gamma(X+Y)$.
In homogeneous degree $n$ it is a scalar multiple of
$X^n+Y^n-(X+Y)^n$. Since multiplication by $XY$ is injective, this bounds the image by the line (D.22). Equivalently the actual module action has leading multiplier $-XYp$, because
$D([x,y])=[x,[p[x,y],y]]=-XYp[x,y]$.
This is the leading term of the full profinite beta multiplier in (D.19), rather than an unlabelled identification of the two series.

The program \texttt{d\_gamma.py} checks $M(h)=XY\bar h$ on all 224 Lyndon basis words of degrees 2 through 10, and solves for each gamma polynomial in the exact invariant image basis of \texttt{d\_metabelian.json}. In weight nine its coordinates in that basis are $(15/2,9/2)$. The first basis vector $X^2Y^2(X+Y)^3$ is independent of the gamma line, verified by an exact rank-two calculation. Thus the comparison with the extra vector is computed in the same coordinate system as D.11.

The special equation forces $X+Y$ to divide $p$, as proved in D.11. At $Y=-X$ the numerator of (D.22) is $-(1+(-1)^n)X^n$, so every positive even-degree scalar must vanish in characteristic zero. Weight one vanishes directly. For each odd weight 3,5,7,9, E.3 supplies a nonzero Galois class for every $\ell$, using \cite[Theorem 8.4, \S8.4, Lemma 8.5 and Remark 8.6]{E_HainMatsumoto}; at $\ell=37$ D.10 also supplies direct finite-field certificates. Since $G_{\mathbb Q}\subset\widehat{\rm GT}$, these lower bounds equal the upper bounds just proved. This proves (D.21) for actual profinite source groups with rational graded target. It does not replace either source by its pro-unipotent automorphism group.
\end{proof}

\begin{statement}{D.22}{Known}{The exact finite-quotient criterion for the remaining GT0 lifts}
Let $N\triangleleft B_3$ have finite index and lie in $PB_3$, and let $\operatorname{GT}(N)$ be the finite set of pairs satisfying the two hexagon equations, the commutator condition on $f$, the unit condition on $\lambda$, and surjectivity onto $B_3/N$. Write $R_{K,N}$ for reduction from a finer such quotient. Then the image of the actual full profinite group is exactly
\[
 \operatorname{Im}\bigl(\widehat{\rm GT}_0\longrightarrow\operatorname{GT}(N)\bigr)
   =\bigcap_{K\triangleleft B_3,\ [B_3:K]<\infty,\ K\subset N}
          R_{K,N}\bigl(\operatorname{GT}(K)\bigr).             \tag{D.23}
\]
This is \cite[Theorem 5.2, Corollary 5.4 and Proposition 5.5]{dDolgushevGuynee2024}. It includes all finite quotients of the braid group, including the nonnilpotent ones. Consequently the finite Lie and CM solutions in D.11 and D.17 specify candidates and concrete finite tests, but their realization in $\widehat{\rm GT}_0$ requires the intersection in (D.23).
\end{statement}
\begin{proof}
One inclusion is reduction of a genuine element. For the other, the fibres over a fixed member of the right side form an inverse system of nonempty finite sets indexed by all finer quotients. Any finite collection is dominated by their intersection, so the system has the finite intersection property. Compactness gives a compatible family. Theorem 5.2 identifies these compatible pairs with the profinite group of hexagon-preserving invertible automorphisms, and Proposition 5.5 identifies that group with Harbater--Schneps' $\widehat{\rm GT}_0$. Surjectivity in each finite quotient supplies the invertibility condition; it cannot be omitted. This proves the equality, including the exact hypothesis that the pro-$\ell$ lift of D.12 alone does not supply.
\end{proof}

\begin{statement}{D.23}{Proved}{The exact coefficient map from the profinite universal measure to every rational Kummer character}
Use the positive interval paths and the Ihara convention of D.20. Write
$c=[x,y]$, $\Lambda=\widehat{\mathbb Z}[[\widehat{\mathbb Z}^2]]$,
and let $s,t$ act on $\widehat F_2'/\widehat F_2''=\Lambda c$ by conjugation by $x,y$.
For $g\in G_{\mathbb Q}$ put $f_g\bmod\widehat F_2''=u_g c$ and
$A_\lambda(s)=(s^\lambda-1)/(s-1)$, where $\lambda=\chi(g)$.
Then the universal measures of D.5 satisfy the exact identities
\[
 F_g=A_\lambda(s)A_\lambda(t)
          \bigl(1+(s-1)(t-1)u_g\bigr),\qquad
 B_g=1+(s-1)(t-1)u_g.                                    \tag{D.24}
\]
For a positive integer $n$, reduce $u_g$ by $t=1$, $s^n=1$, and write
$u_{g,n}=\sum_{j=0}^{n-1}a_j(g)s^j\in\widehat{\mathbb Z}[\mathbb Z/n]$.
The coefficient map is
\[
 \kappa_n(g)=a_0(g)-a_{n-1}(g).                            \tag{D.25}
\]
Consequently the homomorphism of scan 157 is recovered from $(\chi,F)$ by
\[
 g\longmapsto
 \left(\chi(g),\frac{\chi(g)-1}{2}
                 -\bigl(a_0(g)-a_1(g)\bigr)\right),        \tag{D.26}
\]
where the two coefficients in (D.26) use $n=2$.
For $p\ne\ell$, the extraction map in (D.25) on the $\mathbb Z_\ell$ coefficient component does not factor through the quotient with only pro-$\ell$ deck groups; the proof gives an explicit element of its kernel on which extraction is 1. For every odd $\ell$, this failure already occurs on actual Galois elements: a pro-$\ell$ tame-inertia subgroup at 2 has trivial full cyclotomic character and trivial normalized pro-$\ell$ action on $\pi_1(\mathbb P^1-\{0,1,\infty\})$, but its $\kappa_{2,\ell}$ image is all of $\mathbb Z_\ell$.
\end{statement}
\begin{proof}
All divisions in (D.24) are well-defined continuous operations on their stated images. The product $(s-1)(t-1)$ is a non-zero-divisor by Otsubo's Corollary 6.2, hence its individual factors are also non-zero-divisors. The finite geometric sums defining $A_m$ for integers $m$ extend by continuity to $\lambda\in\widehat{\mathbb Z}$. For a unit $\lambda$,
$A_\lambda(s)A_{\lambda^{-1}}(s^\lambda)=1$, by multiplication by $s-1$ and cancellation, so $A_\lambda$ is a unit. The inverse on a closed divisible image is continuous because the completed group ring is compact and Hausdorff.

In additive notation for the derived module,
$f_g^{-1}y^\lambda f_g=(t^\lambda-1)u_gc\cdot y^\lambda$.
The commutator identities give
\[
 [x^\lambda,y^\lambda]=A_\lambda(s)A_\lambda(t)c,
 \qquad
 [x^\lambda,(t^\lambda-1)u_gc]=(s^\lambda-1)(t^\lambda-1)u_gc.
\]
This proves the first identity of (D.24); Otsubo's comparison (D.5) proves the second. To fix the path convention, the lift of $[x,y]$ along the positive interval has homology chain
$(s-1)(1-t)\delta=-(1-s)(1-t)\delta$. Thus it is the negative of Otsubo's Pochhammer generator. The factor is the constant $-1$, so their Galois multipliers agree for every $\chi$, not only on the cyclotomic kernel. The leading derivation convention in D.11 uses $D(y)=[f,y]$, so its $f$ is the negative of the leading term of this Ihara conjugator; this accounts for the sign $-XYp$ in D.21.

Consider the Schreier generators $e_j=x^{-j}yx^j$ of $H_n$, besides $x^n$. In $H_n^{\rm ab}$,
$c=e_{n-1}-e_0$. Conjugation by $x^a$ sends this to
$e_{-a-1}-e_{-a}$, with indices modulo $n$; conjugation by $y$ acts trivially on this abelianization. Therefore the coefficient of $y$ in $\phi_n(u_gc)$ is $a_{n-1}-a_0$.
By Ihara's exact equality (D.18), that coefficient is $-\kappa_n(g)$, proving (D.25). Formula (D.26) now follows from the chosen roots
$(-1/2)^{1/m}=\zeta_{2m}/2^{1/m}$ and
$\kappa_{-1}=(\chi-1)/2$. The formulas give a continuous map from the actual metabelian action to the affine image of D.1, with no assertion that arbitrary coefficients satisfy Galois relations.

Here is the promised prime-to-$\ell$ kernel element. In
$\mathbb Z_\ell[[\widehat{\mathbb Z}^2]]$, take the normalized Haar idempotent $e_p$ of the $\mathbb Z_p$ factor of the first deck group. Its images at finite $p^r$-levels are
$p^{-r}\sum_{j=0}^{p^r-1}s^j$; these are compatible since $p$ is a unit in $\mathbb Z_\ell$. Projection to the pro-$\ell$ deck quotient sends $e_p$ to 1. At $t=1,s^p=1$ its image is $p^{-1}\sum_{j=0}^{p-1}s^j$. Hence $1-e_p$ lies in the kernel of the pro-$\ell$ deck projection but its coefficient extraction is 1. This proves the exact failure of that factorization on the ambient module.

For the asserted arithmetic witness, take an odd $\ell$ and a pro-$\ell$ Sylow subgroup $J$ of inertia at 2. Its tame quotient is $\mathbb Z_\ell(1)$. Every root of unity of odd order lies in the maximal unramified extension of $\mathbb Q_2$, and the action on the $2$-power roots has pro-$2$ image. Therefore $\chi(J)=1$. The extensions obtained by adjoining $2^{1/\ell^r}$ over the maximal unramified field are totally tamely ramified of degree $\ell^r$; thus $\kappa_{2,\ell}(J)=\mathbb Z_\ell$. The pro-$\ell$ outer action of the thrice-punctured line is unramified at 2, as in E.3 and \cite[\S8.4]{E_HainMatsumoto}. It follows that the normalized Ihara action of any $g\in J$ is inner. An inner automorphism fixing $x$ has conjugator in the centralizer $\langle x\rangle$. Comparing its action on $y$ with $f_g^{-1}yf_g$, where $f_g\in F_\ell'$, shows that $f_g$ differs from a power of $x$ by a power of $y$. Abelianization forces both powers to vanish, so this normalized action is the identity. Hence its pure pro-$\ell$ universal measure is 1, while the full profinite coefficient extraction ranges over $\mathbb Z_\ell$. This proves the asserted failure on actual Galois elements, even when the whole cyclotomic character is retained.

The program \texttt{d\_kummer\_measure.py} independently computes abelian Fox derivatives and Schreier coefficients on 160 explicit free words at all 24 levels $1\leq n\leq24$, giving3840 exact checks of (D.25)'s coefficient identity. It also checks210 commutator-action identities for positive and negative integer exponents, which extend continuously as above. Its output does not substitute computation for the Galois theorem (D.18); that theorem supplies the arithmetic meaning of the checked coefficients.
\end{proof}
